\documentclass[11pt,a4paper]{article}
% ============================================
% GIFT v3.4 — Shared Preamble
% ============================================
% Usage: % ============================================
% GIFT v3.4 — Shared Preamble
% ============================================
% Usage: % ============================================
% GIFT v3.4 — Shared Preamble
% ============================================
% Usage: \input{gift_preamble} after \documentclass

% ENCODING & FONTS
\usepackage[utf8]{inputenc}
\usepackage[T1]{fontenc}
\usepackage{lmodern}

% PAGE LAYOUT (Golden Ratio)
\usepackage[margin=1.618cm, top=2.618cm, bottom=2.618cm]{geometry}

% ESSENTIAL PACKAGES
\usepackage{float}
\usepackage{caption}
\usepackage{subcaption}
\usepackage{setspace}
\usepackage{fancyhdr}
\usepackage{xcolor}
\usepackage{hyperref}
\usepackage{csquotes}
\usepackage{amsmath}
\usepackage{amssymb}
\usepackage{amsthm}
\usepackage{mathtools}
\usepackage{booktabs}
\usepackage{longtable}
\usepackage{array}
\usepackage{multirow}
\usepackage{tikz}
\usepackage{graphicx}
\usepackage{listings}
\usepackage{enumitem}
% Unicode fallbacks (pandoc artifacts that survive markdown→LaTeX)
\DeclareUnicodeCharacter{00B0}{\ensuremath{^\circ}}
\DeclareUnicodeCharacter{221A}{\ensuremath{\sqrt{\,}}}
\DeclareUnicodeCharacter{2713}{\ensuremath{\checkmark}}
\DeclareUnicodeCharacter{2717}{\ensuremath{\times}}
\DeclareUnicodeCharacter{2070}{\ensuremath{^{0}}}
\DeclareUnicodeCharacter{00B9}{\ensuremath{^{1}}}
\DeclareUnicodeCharacter{00B2}{\ensuremath{^{2}}}
\DeclareUnicodeCharacter{00B3}{\ensuremath{^{3}}}
\DeclareUnicodeCharacter{2074}{\ensuremath{^{4}}}
\DeclareUnicodeCharacter{2075}{\ensuremath{^{5}}}
\DeclareUnicodeCharacter{2076}{\ensuremath{^{6}}}
\DeclareUnicodeCharacter{2077}{\ensuremath{^{7}}}
\DeclareUnicodeCharacter{2078}{\ensuremath{^{8}}}
\DeclareUnicodeCharacter{2079}{\ensuremath{^{9}}}
\DeclareUnicodeCharacter{2071}{\ensuremath{^{i}}}
\DeclareUnicodeCharacter{2080}{\ensuremath{_{0}}}
\DeclareUnicodeCharacter{2081}{\ensuremath{_{1}}}
\DeclareUnicodeCharacter{2082}{\ensuremath{_{2}}}
\DeclareUnicodeCharacter{2083}{\ensuremath{_{3}}}
\DeclareUnicodeCharacter{2084}{\ensuremath{_{4}}}
\DeclareUnicodeCharacter{2085}{\ensuremath{_{5}}}
\DeclareUnicodeCharacter{2086}{\ensuremath{_{6}}}
\DeclareUnicodeCharacter{2087}{\ensuremath{_{7}}}
\DeclareUnicodeCharacter{2088}{\ensuremath{_{8}}}
\DeclareUnicodeCharacter{2089}{\ensuremath{_{9}}}
\DeclareUnicodeCharacter{208A}{\ensuremath{_{+}}}
\DeclareUnicodeCharacter{208B}{\ensuremath{_{-}}}
\DeclareUnicodeCharacter{2192}{\ensuremath{\to}}
\DeclareUnicodeCharacter{2502}{|}
\DeclareUnicodeCharacter{25BC}{\ensuremath{\blacktriangledown}}
\DeclareUnicodeCharacter{2500}{\ensuremath{-}}
\DeclareUnicodeCharacter{2514}{\ensuremath{\llcorner}}
\DeclareUnicodeCharacter{251C}{\ensuremath{\vdash}}
\DeclareUnicodeCharacter{2026}{\ldots}
\DeclareUnicodeCharacter{2208}{\ensuremath{\in}}
\DeclareUnicodeCharacter{2229}{\ensuremath{\cap}}
\DeclareUnicodeCharacter{222A}{\ensuremath{\cup}}
\DeclareUnicodeCharacter{2295}{\ensuremath{\oplus}}
\DeclareUnicodeCharacter{2297}{\ensuremath{\otimes}}
\DeclareUnicodeCharacter{2245}{\ensuremath{\cong}}
\DeclareUnicodeCharacter{2248}{\ensuremath{\approx}}
\DeclareUnicodeCharacter{2260}{\ensuremath{\neq}}
\DeclareUnicodeCharacter{2261}{\ensuremath{\equiv}}
\DeclareUnicodeCharacter{2264}{\ensuremath{\leq}}
\DeclareUnicodeCharacter{2265}{\ensuremath{\geq}}
\DeclareUnicodeCharacter{00D7}{\ensuremath{\times}}
\DeclareUnicodeCharacter{00B1}{\ensuremath{\pm}}
\DeclareUnicodeCharacter{2225}{\ensuremath{\|}}
% Greek letters used in body text (pandoc passes them through unchanged)
\DeclareUnicodeCharacter{03B1}{\ensuremath{\alpha}}
\DeclareUnicodeCharacter{03B2}{\ensuremath{\beta}}
\DeclareUnicodeCharacter{03B3}{\ensuremath{\gamma}}
\DeclareUnicodeCharacter{03B4}{\ensuremath{\delta}}
\DeclareUnicodeCharacter{03B5}{\ensuremath{\varepsilon}}
\DeclareUnicodeCharacter{03B6}{\ensuremath{\zeta}}
\DeclareUnicodeCharacter{03B7}{\ensuremath{\eta}}
\DeclareUnicodeCharacter{03B8}{\ensuremath{\theta}}
\DeclareUnicodeCharacter{03B9}{\ensuremath{\iota}}
\DeclareUnicodeCharacter{03BA}{\ensuremath{\kappa}}
\DeclareUnicodeCharacter{03BB}{\ensuremath{\lambda}}
\DeclareUnicodeCharacter{03BC}{\ensuremath{\mu}}
\DeclareUnicodeCharacter{03BD}{\ensuremath{\nu}}
\DeclareUnicodeCharacter{03BE}{\ensuremath{\xi}}
\DeclareUnicodeCharacter{03C0}{\ensuremath{\pi}}
\DeclareUnicodeCharacter{03C1}{\ensuremath{\rho}}
\DeclareUnicodeCharacter{03C3}{\ensuremath{\sigma}}
\DeclareUnicodeCharacter{03C4}{\ensuremath{\tau}}
\DeclareUnicodeCharacter{03C5}{\ensuremath{\upsilon}}
\DeclareUnicodeCharacter{03C6}{\ensuremath{\varphi}}
\DeclareUnicodeCharacter{03C7}{\ensuremath{\chi}}
\DeclareUnicodeCharacter{03C8}{\ensuremath{\psi}}
\DeclareUnicodeCharacter{03C9}{\ensuremath{\omega}}
\DeclareUnicodeCharacter{0394}{\ensuremath{\Delta}}
\DeclareUnicodeCharacter{0398}{\ensuremath{\Theta}}
\DeclareUnicodeCharacter{039B}{\ensuremath{\Lambda}}
\DeclareUnicodeCharacter{03A0}{\ensuremath{\Pi}}
\DeclareUnicodeCharacter{03A3}{\ensuremath{\Sigma}}
\DeclareUnicodeCharacter{03A6}{\ensuremath{\Phi}}
\DeclareUnicodeCharacter{03A8}{\ensuremath{\Psi}}
\DeclareUnicodeCharacter{03A9}{\ensuremath{\Omega}}
\DeclareUnicodeCharacter{2200}{\ensuremath{\forall}}
\DeclareUnicodeCharacter{2203}{\ensuremath{\exists}}
\DeclareUnicodeCharacter{2207}{\ensuremath{\nabla}}
\DeclareUnicodeCharacter{2202}{\ensuremath{\partial}}
\DeclareUnicodeCharacter{2218}{\ensuremath{\circ}}
\DeclareUnicodeCharacter{2032}{\ensuremath{\prime}}
\DeclareUnicodeCharacter{2113}{\ensuremath{\ell}}
\DeclareUnicodeCharacter{210F}{\ensuremath{\hbar}}
\DeclareUnicodeCharacter{27E8}{\ensuremath{\langle}}
\DeclareUnicodeCharacter{27E9}{\ensuremath{\rangle}}
\DeclareUnicodeCharacter{2308}{\ensuremath{\lceil}}
\DeclareUnicodeCharacter{2309}{\ensuremath{\rceil}}
\DeclareUnicodeCharacter{230A}{\ensuremath{\lfloor}}
\DeclareUnicodeCharacter{230B}{\ensuremath{\rfloor}}
\DeclareUnicodeCharacter{222B}{\ensuremath{\int}}
\DeclareUnicodeCharacter{2211}{\ensuremath{\sum}}
\DeclareUnicodeCharacter{220F}{\ensuremath{\prod}}
\DeclareUnicodeCharacter{221E}{\ensuremath{\infty}}
\DeclareUnicodeCharacter{00B5}{\ensuremath{\mu}}
\DeclareUnicodeCharacter{2236}{\ensuremath{\colon}}
\DeclareUnicodeCharacter{2192}{\ensuremath{\to}}
\DeclareUnicodeCharacter{21A6}{\ensuremath{\mapsto}}
\DeclareUnicodeCharacter{21D2}{\ensuremath{\Rightarrow}}
\DeclareUnicodeCharacter{21D4}{\ensuremath{\Leftrightarrow}}
\DeclareUnicodeCharacter{2282}{\ensuremath{\subset}}
\DeclareUnicodeCharacter{2286}{\ensuremath{\subseteq}}
\DeclareUnicodeCharacter{2287}{\ensuremath{\supseteq}}
\DeclareUnicodeCharacter{222B}{\ensuremath{\int}}
\DeclareUnicodeCharacter{2205}{\ensuremath{\emptyset}}
\DeclareUnicodeCharacter{2299}{\ensuremath{\odot}}
\DeclareUnicodeCharacter{2220}{\ensuremath{\angle}}
\DeclareUnicodeCharacter{2135}{\ensuremath{\aleph}}
\DeclareUnicodeCharacter{207B}{\ensuremath{^{-}}}
\DeclareUnicodeCharacter{207A}{\ensuremath{^{+}}}
\DeclareUnicodeCharacter{220E}{\ensuremath{\blacksquare}}
\DeclareUnicodeCharacter{2061}{}
\DeclareUnicodeCharacter{2062}{}
\DeclareUnicodeCharacter{2063}{}
\DeclareUnicodeCharacter{2009}{\,}
\DeclareUnicodeCharacter{200A}{\,}
\DeclareUnicodeCharacter{200B}{}
\DeclareUnicodeCharacter{2212}{\ensuremath{-}}
\DeclareUnicodeCharacter{2194}{\ensuremath{\leftrightarrow}}
\DeclareUnicodeCharacter{2099}{\ensuremath{_{n}}}
\DeclareUnicodeCharacter{2098}{\ensuremath{_{m}}}
\DeclareUnicodeCharacter{2096}{\ensuremath{_{k}}}
\DeclareUnicodeCharacter{2095}{\ensuremath{_{h}}}
\DeclareUnicodeCharacter{2090}{\ensuremath{_{a}}}
\DeclareUnicodeCharacter{2091}{\ensuremath{_{e}}}
\DeclareUnicodeCharacter{2093}{\ensuremath{_{x}}}
\DeclareUnicodeCharacter{1D62}{\ensuremath{_{i}}}
\DeclareUnicodeCharacter{2C7C}{\ensuremath{_{j}}}
\DeclareUnicodeCharacter{2C7B}{\ensuremath{^{e}}}
\DeclareUnicodeCharacter{1D52}{\ensuremath{^{o}}}
\DeclareUnicodeCharacter{1D5}{\ensuremath{^{u}}}
\DeclareUnicodeCharacter{02E1}{\ensuremath{^{l}}}
\DeclareUnicodeCharacter{2295}{\ensuremath{\oplus}}
\DeclareUnicodeCharacter{27FA}{\ensuremath{\Longleftrightarrow}}
\DeclareUnicodeCharacter{27F9}{\ensuremath{\Longrightarrow}}
\DeclareUnicodeCharacter{2299}{\ensuremath{\odot}}
% Provide \tightlist used by pandoc-generated lists (no-op)
\providecommand{\tightlist}{\setlength{\itemsep}{0pt}\setlength{\parskip}{0pt}}
% pandoc longtable column specs: provide \real as passthrough and a
% `none` counter so that calc's \real{x} expansion path works.
\newcounter{none}
\providecommand{\real}[1]{#1}

% LISTINGS
\lstset{
    basicstyle=\small\ttfamily,
    breaklines=true, frame=single,
    keepspaces=true, showstringspaces=false,
    aboveskip=0.8em, belowskip=0.8em
}

% HEADER/FOOTER
\setlength{\headheight}{14pt}
\pagestyle{fancy}
\fancyhf{}
\fancyhead[R]{\thepage}
\renewcommand{\headrulewidth}{0.2pt}

% HYPERREF
\hypersetup{
    colorlinks=true,
    linkcolor=blue, citecolor=blue, urlcolor=blue,
    pdfauthor={Brieuc de La Fourniere}
}

% SPACING
\setstretch{1.2}
\setlength{\parskip}{0.4em}
\setlength{\parindent}{0pt}

% TITLE FORMATTING
\usepackage{titling}
\pretitle{\LARGE\bfseries}
\posttitle{\vspace{-0.4em}}
\preauthor{}
\postauthor{}
\predate{}
\postdate{}
\setlength{\droptitle}{-2.0em}

% CUSTOM COMMANDS — GIFT notation
\newcommand{\E}{\mathrm{E}}
\newcommand{\Gtwo}{\mathrm{G}_2}
\newcommand{\Kseven}{K_7}
\newcommand{\AdS}{\mathrm{AdS}}
\newcommand{\SU}{\mathrm{SU}}
\newcommand{\SO}{\mathrm{SO}}
\newcommand{\U}{\mathrm{U}}
\newcommand{\Sp}{\mathrm{Sp}}
\newcommand{\PSL}{\mathrm{PSL}}
\newcommand{\SM}{\mathrm{SM}}
\newcommand{\Tr}{\mathrm{Tr}}
\newcommand{\rank}{\mathrm{rank}}
\newcommand{\Vol}{\mathrm{Vol}}
\newcommand{\Hol}{\mathrm{Hol}}
\newcommand{\Ric}{\mathrm{Ric}}
\newcommand{\Rm}{\mathrm{Rm}}
\newcommand{\Scal}{\mathrm{Scal}}
\DeclareMathOperator{\diag}{diag}
\DeclareMathOperator{\cond}{cond}
\newcommand{\GIFT}{\textrm{GIFT}}
\newcommand{\CP}{\mathrm{CP}}
\newcommand{\CKM}{\mathrm{CKM}}
\newcommand{\PMNS}{\mathrm{PMNS}}
\newcommand{\EW}{\mathrm{EW}}
\newcommand{\Pl}{\mathrm{Pl}}
\newcommand{\DE}{\mathrm{DE}}
\newcommand{\DM}{\mathrm{DM}}
\newcommand{\KK}{\mathrm{KK}}
\newcommand{\NK}{\mathrm{NK}}
\newcommand{\TCS}{\mathrm{TCS}}
\newcommand{\PINN}{\mathrm{PINN}}
\newcommand{\btwo}{b_2}
\newcommand{\bthree}{b_3}
\newcommand{\typeI}{\textsc{I}}
\newcommand{\typeII}{\textsc{II}}
\newcommand{\typeIII}{\textsc{III}}
\newcommand{\typeIV}{\textsc{IV}}
\newcommand{\proven}{\textsc{Proven}}

% PDF bookmark safety
\pdfstringdefDisableCommands{%
  \def\textsubscript#1{#1}%
  \def\textsuperscript#1{#1}%
  \def\CP{CP}%
  \def\Gtwo{G2}%
  \def\Kseven{K7}%
  \def\GIFT{GIFT}%
  \def\E{E}%
  \def\SM{SM}%
}

% THEOREM ENVIRONMENTS
\newtheorem{theorem}{Theorem}[section]
\newtheorem{proposition}[theorem]{Proposition}
\newtheorem{lemma}[theorem]{Lemma}
\newtheorem{corollary}[theorem]{Corollary}
\theoremstyle{definition}
\newtheorem{definition}[theorem]{Definition}
\theoremstyle{remark}
\newtheorem{remark}[theorem]{Remark}
 after \documentclass

% ENCODING & FONTS
\usepackage[utf8]{inputenc}
\usepackage[T1]{fontenc}
\usepackage{lmodern}

% PAGE LAYOUT (Golden Ratio)
\usepackage[margin=1.618cm, top=2.618cm, bottom=2.618cm]{geometry}

% ESSENTIAL PACKAGES
\usepackage{float}
\usepackage{caption}
\usepackage{subcaption}
\usepackage{setspace}
\usepackage{fancyhdr}
\usepackage{xcolor}
\usepackage{hyperref}
\usepackage{csquotes}
\usepackage{amsmath}
\usepackage{amssymb}
\usepackage{amsthm}
\usepackage{mathtools}
\usepackage{booktabs}
\usepackage{longtable}
\usepackage{array}
\usepackage{multirow}
\usepackage{tikz}
\usepackage{graphicx}
\usepackage{listings}
\usepackage{enumitem}
% Unicode fallbacks (pandoc artifacts that survive markdown→LaTeX)
\DeclareUnicodeCharacter{00B0}{\ensuremath{^\circ}}
\DeclareUnicodeCharacter{221A}{\ensuremath{\sqrt{\,}}}
\DeclareUnicodeCharacter{2713}{\ensuremath{\checkmark}}
\DeclareUnicodeCharacter{2717}{\ensuremath{\times}}
\DeclareUnicodeCharacter{2070}{\ensuremath{^{0}}}
\DeclareUnicodeCharacter{00B9}{\ensuremath{^{1}}}
\DeclareUnicodeCharacter{00B2}{\ensuremath{^{2}}}
\DeclareUnicodeCharacter{00B3}{\ensuremath{^{3}}}
\DeclareUnicodeCharacter{2074}{\ensuremath{^{4}}}
\DeclareUnicodeCharacter{2075}{\ensuremath{^{5}}}
\DeclareUnicodeCharacter{2076}{\ensuremath{^{6}}}
\DeclareUnicodeCharacter{2077}{\ensuremath{^{7}}}
\DeclareUnicodeCharacter{2078}{\ensuremath{^{8}}}
\DeclareUnicodeCharacter{2079}{\ensuremath{^{9}}}
\DeclareUnicodeCharacter{2071}{\ensuremath{^{i}}}
\DeclareUnicodeCharacter{2080}{\ensuremath{_{0}}}
\DeclareUnicodeCharacter{2081}{\ensuremath{_{1}}}
\DeclareUnicodeCharacter{2082}{\ensuremath{_{2}}}
\DeclareUnicodeCharacter{2083}{\ensuremath{_{3}}}
\DeclareUnicodeCharacter{2084}{\ensuremath{_{4}}}
\DeclareUnicodeCharacter{2085}{\ensuremath{_{5}}}
\DeclareUnicodeCharacter{2086}{\ensuremath{_{6}}}
\DeclareUnicodeCharacter{2087}{\ensuremath{_{7}}}
\DeclareUnicodeCharacter{2088}{\ensuremath{_{8}}}
\DeclareUnicodeCharacter{2089}{\ensuremath{_{9}}}
\DeclareUnicodeCharacter{208A}{\ensuremath{_{+}}}
\DeclareUnicodeCharacter{208B}{\ensuremath{_{-}}}
\DeclareUnicodeCharacter{2192}{\ensuremath{\to}}
\DeclareUnicodeCharacter{2502}{|}
\DeclareUnicodeCharacter{25BC}{\ensuremath{\blacktriangledown}}
\DeclareUnicodeCharacter{2500}{\ensuremath{-}}
\DeclareUnicodeCharacter{2514}{\ensuremath{\llcorner}}
\DeclareUnicodeCharacter{251C}{\ensuremath{\vdash}}
\DeclareUnicodeCharacter{2026}{\ldots}
\DeclareUnicodeCharacter{2208}{\ensuremath{\in}}
\DeclareUnicodeCharacter{2229}{\ensuremath{\cap}}
\DeclareUnicodeCharacter{222A}{\ensuremath{\cup}}
\DeclareUnicodeCharacter{2295}{\ensuremath{\oplus}}
\DeclareUnicodeCharacter{2297}{\ensuremath{\otimes}}
\DeclareUnicodeCharacter{2245}{\ensuremath{\cong}}
\DeclareUnicodeCharacter{2248}{\ensuremath{\approx}}
\DeclareUnicodeCharacter{2260}{\ensuremath{\neq}}
\DeclareUnicodeCharacter{2261}{\ensuremath{\equiv}}
\DeclareUnicodeCharacter{2264}{\ensuremath{\leq}}
\DeclareUnicodeCharacter{2265}{\ensuremath{\geq}}
\DeclareUnicodeCharacter{00D7}{\ensuremath{\times}}
\DeclareUnicodeCharacter{00B1}{\ensuremath{\pm}}
\DeclareUnicodeCharacter{2225}{\ensuremath{\|}}
% Greek letters used in body text (pandoc passes them through unchanged)
\DeclareUnicodeCharacter{03B1}{\ensuremath{\alpha}}
\DeclareUnicodeCharacter{03B2}{\ensuremath{\beta}}
\DeclareUnicodeCharacter{03B3}{\ensuremath{\gamma}}
\DeclareUnicodeCharacter{03B4}{\ensuremath{\delta}}
\DeclareUnicodeCharacter{03B5}{\ensuremath{\varepsilon}}
\DeclareUnicodeCharacter{03B6}{\ensuremath{\zeta}}
\DeclareUnicodeCharacter{03B7}{\ensuremath{\eta}}
\DeclareUnicodeCharacter{03B8}{\ensuremath{\theta}}
\DeclareUnicodeCharacter{03B9}{\ensuremath{\iota}}
\DeclareUnicodeCharacter{03BA}{\ensuremath{\kappa}}
\DeclareUnicodeCharacter{03BB}{\ensuremath{\lambda}}
\DeclareUnicodeCharacter{03BC}{\ensuremath{\mu}}
\DeclareUnicodeCharacter{03BD}{\ensuremath{\nu}}
\DeclareUnicodeCharacter{03BE}{\ensuremath{\xi}}
\DeclareUnicodeCharacter{03C0}{\ensuremath{\pi}}
\DeclareUnicodeCharacter{03C1}{\ensuremath{\rho}}
\DeclareUnicodeCharacter{03C3}{\ensuremath{\sigma}}
\DeclareUnicodeCharacter{03C4}{\ensuremath{\tau}}
\DeclareUnicodeCharacter{03C5}{\ensuremath{\upsilon}}
\DeclareUnicodeCharacter{03C6}{\ensuremath{\varphi}}
\DeclareUnicodeCharacter{03C7}{\ensuremath{\chi}}
\DeclareUnicodeCharacter{03C8}{\ensuremath{\psi}}
\DeclareUnicodeCharacter{03C9}{\ensuremath{\omega}}
\DeclareUnicodeCharacter{0394}{\ensuremath{\Delta}}
\DeclareUnicodeCharacter{0398}{\ensuremath{\Theta}}
\DeclareUnicodeCharacter{039B}{\ensuremath{\Lambda}}
\DeclareUnicodeCharacter{03A0}{\ensuremath{\Pi}}
\DeclareUnicodeCharacter{03A3}{\ensuremath{\Sigma}}
\DeclareUnicodeCharacter{03A6}{\ensuremath{\Phi}}
\DeclareUnicodeCharacter{03A8}{\ensuremath{\Psi}}
\DeclareUnicodeCharacter{03A9}{\ensuremath{\Omega}}
\DeclareUnicodeCharacter{2200}{\ensuremath{\forall}}
\DeclareUnicodeCharacter{2203}{\ensuremath{\exists}}
\DeclareUnicodeCharacter{2207}{\ensuremath{\nabla}}
\DeclareUnicodeCharacter{2202}{\ensuremath{\partial}}
\DeclareUnicodeCharacter{2218}{\ensuremath{\circ}}
\DeclareUnicodeCharacter{2032}{\ensuremath{\prime}}
\DeclareUnicodeCharacter{2113}{\ensuremath{\ell}}
\DeclareUnicodeCharacter{210F}{\ensuremath{\hbar}}
\DeclareUnicodeCharacter{27E8}{\ensuremath{\langle}}
\DeclareUnicodeCharacter{27E9}{\ensuremath{\rangle}}
\DeclareUnicodeCharacter{2308}{\ensuremath{\lceil}}
\DeclareUnicodeCharacter{2309}{\ensuremath{\rceil}}
\DeclareUnicodeCharacter{230A}{\ensuremath{\lfloor}}
\DeclareUnicodeCharacter{230B}{\ensuremath{\rfloor}}
\DeclareUnicodeCharacter{222B}{\ensuremath{\int}}
\DeclareUnicodeCharacter{2211}{\ensuremath{\sum}}
\DeclareUnicodeCharacter{220F}{\ensuremath{\prod}}
\DeclareUnicodeCharacter{221E}{\ensuremath{\infty}}
\DeclareUnicodeCharacter{00B5}{\ensuremath{\mu}}
\DeclareUnicodeCharacter{2236}{\ensuremath{\colon}}
\DeclareUnicodeCharacter{2192}{\ensuremath{\to}}
\DeclareUnicodeCharacter{21A6}{\ensuremath{\mapsto}}
\DeclareUnicodeCharacter{21D2}{\ensuremath{\Rightarrow}}
\DeclareUnicodeCharacter{21D4}{\ensuremath{\Leftrightarrow}}
\DeclareUnicodeCharacter{2282}{\ensuremath{\subset}}
\DeclareUnicodeCharacter{2286}{\ensuremath{\subseteq}}
\DeclareUnicodeCharacter{2287}{\ensuremath{\supseteq}}
\DeclareUnicodeCharacter{222B}{\ensuremath{\int}}
\DeclareUnicodeCharacter{2205}{\ensuremath{\emptyset}}
\DeclareUnicodeCharacter{2299}{\ensuremath{\odot}}
\DeclareUnicodeCharacter{2220}{\ensuremath{\angle}}
\DeclareUnicodeCharacter{2135}{\ensuremath{\aleph}}
\DeclareUnicodeCharacter{207B}{\ensuremath{^{-}}}
\DeclareUnicodeCharacter{207A}{\ensuremath{^{+}}}
\DeclareUnicodeCharacter{220E}{\ensuremath{\blacksquare}}
\DeclareUnicodeCharacter{2061}{}
\DeclareUnicodeCharacter{2062}{}
\DeclareUnicodeCharacter{2063}{}
\DeclareUnicodeCharacter{2009}{\,}
\DeclareUnicodeCharacter{200A}{\,}
\DeclareUnicodeCharacter{200B}{}
\DeclareUnicodeCharacter{2212}{\ensuremath{-}}
\DeclareUnicodeCharacter{2194}{\ensuremath{\leftrightarrow}}
\DeclareUnicodeCharacter{2099}{\ensuremath{_{n}}}
\DeclareUnicodeCharacter{2098}{\ensuremath{_{m}}}
\DeclareUnicodeCharacter{2096}{\ensuremath{_{k}}}
\DeclareUnicodeCharacter{2095}{\ensuremath{_{h}}}
\DeclareUnicodeCharacter{2090}{\ensuremath{_{a}}}
\DeclareUnicodeCharacter{2091}{\ensuremath{_{e}}}
\DeclareUnicodeCharacter{2093}{\ensuremath{_{x}}}
\DeclareUnicodeCharacter{1D62}{\ensuremath{_{i}}}
\DeclareUnicodeCharacter{2C7C}{\ensuremath{_{j}}}
\DeclareUnicodeCharacter{2C7B}{\ensuremath{^{e}}}
\DeclareUnicodeCharacter{1D52}{\ensuremath{^{o}}}
\DeclareUnicodeCharacter{1D5}{\ensuremath{^{u}}}
\DeclareUnicodeCharacter{02E1}{\ensuremath{^{l}}}
\DeclareUnicodeCharacter{2295}{\ensuremath{\oplus}}
\DeclareUnicodeCharacter{27FA}{\ensuremath{\Longleftrightarrow}}
\DeclareUnicodeCharacter{27F9}{\ensuremath{\Longrightarrow}}
\DeclareUnicodeCharacter{2299}{\ensuremath{\odot}}
% Provide \tightlist used by pandoc-generated lists (no-op)
\providecommand{\tightlist}{\setlength{\itemsep}{0pt}\setlength{\parskip}{0pt}}
% pandoc longtable column specs: provide \real as passthrough and a
% `none` counter so that calc's \real{x} expansion path works.
\newcounter{none}
\providecommand{\real}[1]{#1}

% LISTINGS
\lstset{
    basicstyle=\small\ttfamily,
    breaklines=true, frame=single,
    keepspaces=true, showstringspaces=false,
    aboveskip=0.8em, belowskip=0.8em
}

% HEADER/FOOTER
\setlength{\headheight}{14pt}
\pagestyle{fancy}
\fancyhf{}
\fancyhead[R]{\thepage}
\renewcommand{\headrulewidth}{0.2pt}

% HYPERREF
\hypersetup{
    colorlinks=true,
    linkcolor=blue, citecolor=blue, urlcolor=blue,
    pdfauthor={Brieuc de La Fourniere}
}

% SPACING
\setstretch{1.2}
\setlength{\parskip}{0.4em}
\setlength{\parindent}{0pt}

% TITLE FORMATTING
\usepackage{titling}
\pretitle{\LARGE\bfseries}
\posttitle{\vspace{-0.4em}}
\preauthor{}
\postauthor{}
\predate{}
\postdate{}
\setlength{\droptitle}{-2.0em}

% CUSTOM COMMANDS — GIFT notation
\newcommand{\E}{\mathrm{E}}
\newcommand{\Gtwo}{\mathrm{G}_2}
\newcommand{\Kseven}{K_7}
\newcommand{\AdS}{\mathrm{AdS}}
\newcommand{\SU}{\mathrm{SU}}
\newcommand{\SO}{\mathrm{SO}}
\newcommand{\U}{\mathrm{U}}
\newcommand{\Sp}{\mathrm{Sp}}
\newcommand{\PSL}{\mathrm{PSL}}
\newcommand{\SM}{\mathrm{SM}}
\newcommand{\Tr}{\mathrm{Tr}}
\newcommand{\rank}{\mathrm{rank}}
\newcommand{\Vol}{\mathrm{Vol}}
\newcommand{\Hol}{\mathrm{Hol}}
\newcommand{\Ric}{\mathrm{Ric}}
\newcommand{\Rm}{\mathrm{Rm}}
\newcommand{\Scal}{\mathrm{Scal}}
\DeclareMathOperator{\diag}{diag}
\DeclareMathOperator{\cond}{cond}
\newcommand{\GIFT}{\textrm{GIFT}}
\newcommand{\CP}{\mathrm{CP}}
\newcommand{\CKM}{\mathrm{CKM}}
\newcommand{\PMNS}{\mathrm{PMNS}}
\newcommand{\EW}{\mathrm{EW}}
\newcommand{\Pl}{\mathrm{Pl}}
\newcommand{\DE}{\mathrm{DE}}
\newcommand{\DM}{\mathrm{DM}}
\newcommand{\KK}{\mathrm{KK}}
\newcommand{\NK}{\mathrm{NK}}
\newcommand{\TCS}{\mathrm{TCS}}
\newcommand{\PINN}{\mathrm{PINN}}
\newcommand{\btwo}{b_2}
\newcommand{\bthree}{b_3}
\newcommand{\typeI}{\textsc{I}}
\newcommand{\typeII}{\textsc{II}}
\newcommand{\typeIII}{\textsc{III}}
\newcommand{\typeIV}{\textsc{IV}}
\newcommand{\proven}{\textsc{Proven}}

% PDF bookmark safety
\pdfstringdefDisableCommands{%
  \def\textsubscript#1{#1}%
  \def\textsuperscript#1{#1}%
  \def\CP{CP}%
  \def\Gtwo{G2}%
  \def\Kseven{K7}%
  \def\GIFT{GIFT}%
  \def\E{E}%
  \def\SM{SM}%
}

% THEOREM ENVIRONMENTS
\newtheorem{theorem}{Theorem}[section]
\newtheorem{proposition}[theorem]{Proposition}
\newtheorem{lemma}[theorem]{Lemma}
\newtheorem{corollary}[theorem]{Corollary}
\theoremstyle{definition}
\newtheorem{definition}[theorem]{Definition}
\theoremstyle{remark}
\newtheorem{remark}[theorem]{Remark}
 after \documentclass

% ENCODING & FONTS
\usepackage[utf8]{inputenc}
\usepackage[T1]{fontenc}
\usepackage{lmodern}

% PAGE LAYOUT (Golden Ratio)
\usepackage[margin=1.618cm, top=2.618cm, bottom=2.618cm]{geometry}

% ESSENTIAL PACKAGES
\usepackage{float}
\usepackage{caption}
\usepackage{subcaption}
\usepackage{setspace}
\usepackage{fancyhdr}
\usepackage{xcolor}
\usepackage{hyperref}
\usepackage{csquotes}
\usepackage{amsmath}
\usepackage{amssymb}
\usepackage{amsthm}
\usepackage{mathtools}
\usepackage{booktabs}
\usepackage{longtable}
\usepackage{array}
\usepackage{multirow}
\usepackage{tikz}
\usepackage{graphicx}
\usepackage{listings}
\usepackage{enumitem}
% Unicode fallbacks (pandoc artifacts that survive markdown→LaTeX)
\DeclareUnicodeCharacter{00B0}{\ensuremath{^\circ}}
\DeclareUnicodeCharacter{221A}{\ensuremath{\sqrt{\,}}}
\DeclareUnicodeCharacter{2713}{\ensuremath{\checkmark}}
\DeclareUnicodeCharacter{2717}{\ensuremath{\times}}
\DeclareUnicodeCharacter{2070}{\ensuremath{^{0}}}
\DeclareUnicodeCharacter{00B9}{\ensuremath{^{1}}}
\DeclareUnicodeCharacter{00B2}{\ensuremath{^{2}}}
\DeclareUnicodeCharacter{00B3}{\ensuremath{^{3}}}
\DeclareUnicodeCharacter{2074}{\ensuremath{^{4}}}
\DeclareUnicodeCharacter{2075}{\ensuremath{^{5}}}
\DeclareUnicodeCharacter{2076}{\ensuremath{^{6}}}
\DeclareUnicodeCharacter{2077}{\ensuremath{^{7}}}
\DeclareUnicodeCharacter{2078}{\ensuremath{^{8}}}
\DeclareUnicodeCharacter{2079}{\ensuremath{^{9}}}
\DeclareUnicodeCharacter{2071}{\ensuremath{^{i}}}
\DeclareUnicodeCharacter{2080}{\ensuremath{_{0}}}
\DeclareUnicodeCharacter{2081}{\ensuremath{_{1}}}
\DeclareUnicodeCharacter{2082}{\ensuremath{_{2}}}
\DeclareUnicodeCharacter{2083}{\ensuremath{_{3}}}
\DeclareUnicodeCharacter{2084}{\ensuremath{_{4}}}
\DeclareUnicodeCharacter{2085}{\ensuremath{_{5}}}
\DeclareUnicodeCharacter{2086}{\ensuremath{_{6}}}
\DeclareUnicodeCharacter{2087}{\ensuremath{_{7}}}
\DeclareUnicodeCharacter{2088}{\ensuremath{_{8}}}
\DeclareUnicodeCharacter{2089}{\ensuremath{_{9}}}
\DeclareUnicodeCharacter{208A}{\ensuremath{_{+}}}
\DeclareUnicodeCharacter{208B}{\ensuremath{_{-}}}
\DeclareUnicodeCharacter{2192}{\ensuremath{\to}}
\DeclareUnicodeCharacter{2502}{|}
\DeclareUnicodeCharacter{25BC}{\ensuremath{\blacktriangledown}}
\DeclareUnicodeCharacter{2500}{\ensuremath{-}}
\DeclareUnicodeCharacter{2514}{\ensuremath{\llcorner}}
\DeclareUnicodeCharacter{251C}{\ensuremath{\vdash}}
\DeclareUnicodeCharacter{2026}{\ldots}
\DeclareUnicodeCharacter{2208}{\ensuremath{\in}}
\DeclareUnicodeCharacter{2229}{\ensuremath{\cap}}
\DeclareUnicodeCharacter{222A}{\ensuremath{\cup}}
\DeclareUnicodeCharacter{2295}{\ensuremath{\oplus}}
\DeclareUnicodeCharacter{2297}{\ensuremath{\otimes}}
\DeclareUnicodeCharacter{2245}{\ensuremath{\cong}}
\DeclareUnicodeCharacter{2248}{\ensuremath{\approx}}
\DeclareUnicodeCharacter{2260}{\ensuremath{\neq}}
\DeclareUnicodeCharacter{2261}{\ensuremath{\equiv}}
\DeclareUnicodeCharacter{2264}{\ensuremath{\leq}}
\DeclareUnicodeCharacter{2265}{\ensuremath{\geq}}
\DeclareUnicodeCharacter{00D7}{\ensuremath{\times}}
\DeclareUnicodeCharacter{00B1}{\ensuremath{\pm}}
\DeclareUnicodeCharacter{2225}{\ensuremath{\|}}
% Greek letters used in body text (pandoc passes them through unchanged)
\DeclareUnicodeCharacter{03B1}{\ensuremath{\alpha}}
\DeclareUnicodeCharacter{03B2}{\ensuremath{\beta}}
\DeclareUnicodeCharacter{03B3}{\ensuremath{\gamma}}
\DeclareUnicodeCharacter{03B4}{\ensuremath{\delta}}
\DeclareUnicodeCharacter{03B5}{\ensuremath{\varepsilon}}
\DeclareUnicodeCharacter{03B6}{\ensuremath{\zeta}}
\DeclareUnicodeCharacter{03B7}{\ensuremath{\eta}}
\DeclareUnicodeCharacter{03B8}{\ensuremath{\theta}}
\DeclareUnicodeCharacter{03B9}{\ensuremath{\iota}}
\DeclareUnicodeCharacter{03BA}{\ensuremath{\kappa}}
\DeclareUnicodeCharacter{03BB}{\ensuremath{\lambda}}
\DeclareUnicodeCharacter{03BC}{\ensuremath{\mu}}
\DeclareUnicodeCharacter{03BD}{\ensuremath{\nu}}
\DeclareUnicodeCharacter{03BE}{\ensuremath{\xi}}
\DeclareUnicodeCharacter{03C0}{\ensuremath{\pi}}
\DeclareUnicodeCharacter{03C1}{\ensuremath{\rho}}
\DeclareUnicodeCharacter{03C3}{\ensuremath{\sigma}}
\DeclareUnicodeCharacter{03C4}{\ensuremath{\tau}}
\DeclareUnicodeCharacter{03C5}{\ensuremath{\upsilon}}
\DeclareUnicodeCharacter{03C6}{\ensuremath{\varphi}}
\DeclareUnicodeCharacter{03C7}{\ensuremath{\chi}}
\DeclareUnicodeCharacter{03C8}{\ensuremath{\psi}}
\DeclareUnicodeCharacter{03C9}{\ensuremath{\omega}}
\DeclareUnicodeCharacter{0394}{\ensuremath{\Delta}}
\DeclareUnicodeCharacter{0398}{\ensuremath{\Theta}}
\DeclareUnicodeCharacter{039B}{\ensuremath{\Lambda}}
\DeclareUnicodeCharacter{03A0}{\ensuremath{\Pi}}
\DeclareUnicodeCharacter{03A3}{\ensuremath{\Sigma}}
\DeclareUnicodeCharacter{03A6}{\ensuremath{\Phi}}
\DeclareUnicodeCharacter{03A8}{\ensuremath{\Psi}}
\DeclareUnicodeCharacter{03A9}{\ensuremath{\Omega}}
\DeclareUnicodeCharacter{2200}{\ensuremath{\forall}}
\DeclareUnicodeCharacter{2203}{\ensuremath{\exists}}
\DeclareUnicodeCharacter{2207}{\ensuremath{\nabla}}
\DeclareUnicodeCharacter{2202}{\ensuremath{\partial}}
\DeclareUnicodeCharacter{2218}{\ensuremath{\circ}}
\DeclareUnicodeCharacter{2032}{\ensuremath{\prime}}
\DeclareUnicodeCharacter{2113}{\ensuremath{\ell}}
\DeclareUnicodeCharacter{210F}{\ensuremath{\hbar}}
\DeclareUnicodeCharacter{27E8}{\ensuremath{\langle}}
\DeclareUnicodeCharacter{27E9}{\ensuremath{\rangle}}
\DeclareUnicodeCharacter{2308}{\ensuremath{\lceil}}
\DeclareUnicodeCharacter{2309}{\ensuremath{\rceil}}
\DeclareUnicodeCharacter{230A}{\ensuremath{\lfloor}}
\DeclareUnicodeCharacter{230B}{\ensuremath{\rfloor}}
\DeclareUnicodeCharacter{222B}{\ensuremath{\int}}
\DeclareUnicodeCharacter{2211}{\ensuremath{\sum}}
\DeclareUnicodeCharacter{220F}{\ensuremath{\prod}}
\DeclareUnicodeCharacter{221E}{\ensuremath{\infty}}
\DeclareUnicodeCharacter{00B5}{\ensuremath{\mu}}
\DeclareUnicodeCharacter{2236}{\ensuremath{\colon}}
\DeclareUnicodeCharacter{2192}{\ensuremath{\to}}
\DeclareUnicodeCharacter{21A6}{\ensuremath{\mapsto}}
\DeclareUnicodeCharacter{21D2}{\ensuremath{\Rightarrow}}
\DeclareUnicodeCharacter{21D4}{\ensuremath{\Leftrightarrow}}
\DeclareUnicodeCharacter{2282}{\ensuremath{\subset}}
\DeclareUnicodeCharacter{2286}{\ensuremath{\subseteq}}
\DeclareUnicodeCharacter{2287}{\ensuremath{\supseteq}}
\DeclareUnicodeCharacter{222B}{\ensuremath{\int}}
\DeclareUnicodeCharacter{2205}{\ensuremath{\emptyset}}
\DeclareUnicodeCharacter{2299}{\ensuremath{\odot}}
\DeclareUnicodeCharacter{2220}{\ensuremath{\angle}}
\DeclareUnicodeCharacter{2135}{\ensuremath{\aleph}}
\DeclareUnicodeCharacter{207B}{\ensuremath{^{-}}}
\DeclareUnicodeCharacter{207A}{\ensuremath{^{+}}}
\DeclareUnicodeCharacter{220E}{\ensuremath{\blacksquare}}
\DeclareUnicodeCharacter{2061}{}
\DeclareUnicodeCharacter{2062}{}
\DeclareUnicodeCharacter{2063}{}
\DeclareUnicodeCharacter{2009}{\,}
\DeclareUnicodeCharacter{200A}{\,}
\DeclareUnicodeCharacter{200B}{}
\DeclareUnicodeCharacter{2212}{\ensuremath{-}}
\DeclareUnicodeCharacter{2194}{\ensuremath{\leftrightarrow}}
\DeclareUnicodeCharacter{2099}{\ensuremath{_{n}}}
\DeclareUnicodeCharacter{2098}{\ensuremath{_{m}}}
\DeclareUnicodeCharacter{2096}{\ensuremath{_{k}}}
\DeclareUnicodeCharacter{2095}{\ensuremath{_{h}}}
\DeclareUnicodeCharacter{2090}{\ensuremath{_{a}}}
\DeclareUnicodeCharacter{2091}{\ensuremath{_{e}}}
\DeclareUnicodeCharacter{2093}{\ensuremath{_{x}}}
\DeclareUnicodeCharacter{1D62}{\ensuremath{_{i}}}
\DeclareUnicodeCharacter{2C7C}{\ensuremath{_{j}}}
\DeclareUnicodeCharacter{2C7B}{\ensuremath{^{e}}}
\DeclareUnicodeCharacter{1D52}{\ensuremath{^{o}}}
\DeclareUnicodeCharacter{1D5}{\ensuremath{^{u}}}
\DeclareUnicodeCharacter{02E1}{\ensuremath{^{l}}}
\DeclareUnicodeCharacter{2295}{\ensuremath{\oplus}}
\DeclareUnicodeCharacter{27FA}{\ensuremath{\Longleftrightarrow}}
\DeclareUnicodeCharacter{27F9}{\ensuremath{\Longrightarrow}}
\DeclareUnicodeCharacter{2299}{\ensuremath{\odot}}
% Provide \tightlist used by pandoc-generated lists (no-op)
\providecommand{\tightlist}{\setlength{\itemsep}{0pt}\setlength{\parskip}{0pt}}
% pandoc longtable column specs: provide \real as passthrough and a
% `none` counter so that calc's \real{x} expansion path works.
\newcounter{none}
\providecommand{\real}[1]{#1}

% LISTINGS
\lstset{
    basicstyle=\small\ttfamily,
    breaklines=true, frame=single,
    keepspaces=true, showstringspaces=false,
    aboveskip=0.8em, belowskip=0.8em
}

% HEADER/FOOTER
\setlength{\headheight}{14pt}
\pagestyle{fancy}
\fancyhf{}
\fancyhead[R]{\thepage}
\renewcommand{\headrulewidth}{0.2pt}

% HYPERREF
\hypersetup{
    colorlinks=true,
    linkcolor=blue, citecolor=blue, urlcolor=blue,
    pdfauthor={Brieuc de La Fourniere}
}

% SPACING
\setstretch{1.2}
\setlength{\parskip}{0.4em}
\setlength{\parindent}{0pt}

% TITLE FORMATTING
\usepackage{titling}
\pretitle{\LARGE\bfseries}
\posttitle{\vspace{-0.4em}}
\preauthor{}
\postauthor{}
\predate{}
\postdate{}
\setlength{\droptitle}{-2.0em}

% CUSTOM COMMANDS — GIFT notation
\newcommand{\E}{\mathrm{E}}
\newcommand{\Gtwo}{\mathrm{G}_2}
\newcommand{\Kseven}{K_7}
\newcommand{\AdS}{\mathrm{AdS}}
\newcommand{\SU}{\mathrm{SU}}
\newcommand{\SO}{\mathrm{SO}}
\newcommand{\U}{\mathrm{U}}
\newcommand{\Sp}{\mathrm{Sp}}
\newcommand{\PSL}{\mathrm{PSL}}
\newcommand{\SM}{\mathrm{SM}}
\newcommand{\Tr}{\mathrm{Tr}}
\newcommand{\rank}{\mathrm{rank}}
\newcommand{\Vol}{\mathrm{Vol}}
\newcommand{\Hol}{\mathrm{Hol}}
\newcommand{\Ric}{\mathrm{Ric}}
\newcommand{\Rm}{\mathrm{Rm}}
\newcommand{\Scal}{\mathrm{Scal}}
\DeclareMathOperator{\diag}{diag}
\DeclareMathOperator{\cond}{cond}
\newcommand{\GIFT}{\textrm{GIFT}}
\newcommand{\CP}{\mathrm{CP}}
\newcommand{\CKM}{\mathrm{CKM}}
\newcommand{\PMNS}{\mathrm{PMNS}}
\newcommand{\EW}{\mathrm{EW}}
\newcommand{\Pl}{\mathrm{Pl}}
\newcommand{\DE}{\mathrm{DE}}
\newcommand{\DM}{\mathrm{DM}}
\newcommand{\KK}{\mathrm{KK}}
\newcommand{\NK}{\mathrm{NK}}
\newcommand{\TCS}{\mathrm{TCS}}
\newcommand{\PINN}{\mathrm{PINN}}
\newcommand{\btwo}{b_2}
\newcommand{\bthree}{b_3}
\newcommand{\typeI}{\textsc{I}}
\newcommand{\typeII}{\textsc{II}}
\newcommand{\typeIII}{\textsc{III}}
\newcommand{\typeIV}{\textsc{IV}}
\newcommand{\proven}{\textsc{Proven}}

% PDF bookmark safety
\pdfstringdefDisableCommands{%
  \def\textsubscript#1{#1}%
  \def\textsuperscript#1{#1}%
  \def\CP{CP}%
  \def\Gtwo{G2}%
  \def\Kseven{K7}%
  \def\GIFT{GIFT}%
  \def\E{E}%
  \def\SM{SM}%
}

% THEOREM ENVIRONMENTS
\newtheorem{theorem}{Theorem}[section]
\newtheorem{proposition}[theorem]{Proposition}
\newtheorem{lemma}[theorem]{Lemma}
\newtheorem{corollary}[theorem]{Corollary}
\theoremstyle{definition}
\newtheorem{definition}[theorem]{Definition}
\theoremstyle{remark}
\newtheorem{remark}[theorem]{Remark}

\fancyhead[L]{GIFT v3.4: Standard Model Parameters as Topological Invariants}
\hypersetup{pdftitle={GIFT v3.4: Standard Model Parameters as Topological Invariants of a G2 Holonomy Manifold}}

\graphicspath{{../../../figures/}}

% ============================================
% TITLE
% ============================================
\title{%
\LARGE\textbf{Geometric Information Field Theory}\\[0.3em]
\large Standard Model Parameters as Topological Invariants\\
of a $\mathrm{G}_2$ Holonomy Manifold%
}
\author{}
\date{}

\begin{document}

\maketitle
\noindent\rule{\textwidth}{0.2pt}

\noindent\textbf{Author}: Brieuc de La Fourni\`ere

\noindent Independent researcher, Beaune, France

\noindent Contact: \href{mailto:brieuc.dlf@gmail.com}{\texttt{brieuc.dlf@gmail.com}}

\noindent\textbf{Date}: May 2026 \quad v3.4

\noindent\textbf{DOI}: \href{https://doi.org/10.5281/zenodo.18837071}{10.5281/zenodo.18837071}

\vfill

\begin{abstract}

The Standard Model requires 19 experimentally determined parameters
lacking theoretical explanation. We present a geometric framework in
which physical observables emerge as topological invariants of a
seven-dimensional \(G_{2}\) holonomy manifold \(K_{7}\) coupled to
\(E_{8} \times E_{8}\) gauge structure, with the topological inputs treated as discrete and the metric determinant det(g) = 65/32 imposed as a normalization target (§2.3); no continuously adjustable physical parameter is introduced.

Building on a Newton--Kantorovich-certified neck/metric model
([A], [B]) and a Joyce--Karigiannis topological/lattice route
realizing \((b_{2}, b_{3}) = (21, 77)\), while the full smooth
analytic compact construction remains open: we derive \textbf{95
observables} from \(K_{7}\) organized in four types: 33 direct
algebraic (Type I, mean deviation 0.73\%), 19 one-step physical
extractions (Type II, 0.17\%), 21 multi-step dynamical chains (Type III,
3.4\%), and 22 structural diagnostics (Type IV). Of 66 experimentally
comparable observables, 11 are exact matches (deviation \textless{}
0.01\%) and 53 are within 1\%.

The framework includes three new framework-level results beyond the original 33
predictions: (1) a proposed \(E_{8}\) \(\to\) Standard Model breaking chain with anomaly-cancellation checks and bundle-universality diagnostics, (2) a
combined lepton mass hierarchy mechanism achieving sub-percent precision
from two independent geometric sources with \(\alpha\) = \(e^K\) (zero
free parameters), and (3) a sensitivity analysis demonstrating
2.13\(\times\) overdetermination with coincidence probability
10\(^{-346}\) under a uniform null and 10\(^{-133}\) under an algebraic
null model (4.2M random formulas from the same 20 constants). Of the 95
observables, 55 are formally verified in Lean 4 (213 certificate
conjuncts, 4 axioms, 0 sorry); Lean certifies the algebraic relations
and internal consistency, not the physical interpretation or the
existence of the complete compact construction (§8.3).

DUNE (2028--2040) will test the topological prediction \(\delta\)\_CP =
197\({}^{\circ}\) with resolution of a few degrees to
\textasciitilde15\({}^{\circ}\); measurement outside {[}182,
212{]}\({}^{\circ}\) would create serious tension. We present this as an
exploratory investigation emphasizing falsifiability, not a claim of
correctness.

\textbf{Keywords}: \(G_{2}\) holonomy, exceptional Lie algebras,
Standard Model parameters, topological field theory, falsifiability,
formal verification, Lean 4

\end{abstract}

\noindent\rule{\textwidth}{0.2pt}

\newpage
\tableofcontents

\bigskip
\noindent\textbf{Framework epistemic status at a glance:}

\begin{center}
\begin{tabular}{lp{9cm}}
\toprule
\textbf{Component} & \textbf{Status} \\
\midrule
Algebraic predictions (33 Type I observables) & Lean-certified (algebraic relations) \\
$K_7$ Betti/lattice gate, $(b_2, b_3) = (21, 77)$ & Certified: NK metric [A] + JK17 topological/lattice route \\
Full smooth compact analytic construction & \textbf{Open} \\
Physical interpretation & Conjectural / falsifiable \\
\bottomrule
\end{tabular}
\end{center}

\newpage

\section{Introduction}

\subsection{The Parameter Problem}

The Standard Model describes fundamental interactions with remarkable
precision, yet requires 19 free parameters determined solely through
experiment \cite{1}. These parameters (gauge couplings, Yukawa couplings
spanning five orders of magnitude, mixing matrices, and Higgs sector
values) lack theoretical explanation.

Several tensions motivate the search for deeper structure:

\begin{itemize}

\item
  \textbf{Hierarchy problem}: The Higgs mass requires fine-tuning absent
  new physics \cite{2}.
\item
  \textbf{Hubble tension}: CMB and local \(H_{0}\) measurements differ
  by \textgreater4 sigma {[}3,4{]}.
\item
  \textbf{Flavor puzzle}: No mechanism explains three generations or
  mass hierarchies \cite{5}.
\item
  \textbf{Koide mystery}: The charged lepton relation Q = 2/3 holds for
  43 years without explanation \cite{6}.
\end{itemize}

These challenges suggest examining whether parameters might emerge from
geometric or topological structures.

\subsection{Framework Overview}

The Geometric Information Field Theory (GIFT) proposes that
dimensionless parameters represent topological invariants of an
eleven-dimensional spacetime:

\[
E_8 \times E_8 \;(\text{496D}) \;\longrightarrow\; \mathrm{AdS}_4 \times K_7 \;(\text{11D}) \;\longrightarrow\; \text{Standard Model}\;(\text{4D})
\]

The key elements:

\begin{enumerate}


\item
  \textbf{\(E_{8}\) x \(E_{8}\) gauge structure} (dimension 496)
\item
  \textbf{Compact 7-manifold \(K_{7}\)} with \(G_{2}\) holonomy
  (\(b_{2}\) = 21, \(b_{3}\) = 77)
\item
  \textbf{Metric normalization constraint}: det(g) = 65/32, structurally
  motivated by \(G_{2}\)/\(E_{8}\) constants (§2.3) but imposed as a
  normalization target, not derived from topology alone
\item
  \textbf{Cohomological mapping}: Betti numbers constrain field content
\end{enumerate}

We emphasize this represents mathematical exploration, not a claim that
nature realizes this structure. The framework's merit lies in
falsifiable predictions from topological inputs.

\subsection{Paper Organization}

\begin{itemize}

\item
  \textbf{Section 2}: Mathematical framework (\(E_{8} \times E_{8}\),
  \(K_{7}\), \(G_{2}\) structure)
\item
  \textbf{Section 3}: Methodology, epistemic status, and type
  classification
\item
  \textbf{Section 4}: 95 observables across 4 types (33(I) + 19(II) +
  21(III) + 22(IV))
\item
  \textbf{Section 5}: Gauge sector, \(E_{8}\) \(\to\) SM breaking,
  anomaly cancellation, B-test identity, bundle universality
\item
  \textbf{Section 6}: Mass hierarchy: Wilson lines, instantons,
  combined wilson\_line+instanton pipeline
\item
  \textbf{Section 7}: Sensitivity analysis, effective DOF,
  cross-correlations, coincidence test
\item
  \textbf{Section 8}: Formal verification (213 Lean conjuncts, 4 axioms)
  and statistical uniqueness
\item
  \textbf{Section 9}: Discussion, falsifiability, and conclusion
\end{itemize}

Three supplements provide technical details: S1 (Mathematical
Foundations), S2 (Complete Derivations), S3 (Observable Dataset with
full 95-entry table).

\noindent

\rule{\textwidth}{0.2pt}

\section{Mathematical Framework}

\subsection{The Octonionic Foundation}

GIFT emerges from the algebraic fact that \textbf{the octonions are the
largest normed division algebra}.

\begin{longtable}[]{@{}llll@{}}
\toprule\noalign{}
Algebra & Dim & Physics Role & Extends? \\
\midrule\noalign{}
\endhead
\bottomrule\noalign{}
\endlastfoot
R & 1 & Classical mechanics & Yes \\
C & 2 & Quantum mechanics & Yes \\
H & 4 & Spin, Lorentz group & Yes \\
\textbf{O} & \textbf{8} & \textbf{Exceptional structures} &
\textbf{No} \\
\end{longtable}


The octonions terminate this sequence. Their automorphism group
\(G_{2}\) = Aut(O) has dimension 14 and acts naturally on Im(O) =
\(R^7\). The exceptional Lie algebras arise from octonionic
constructions through a chain established by Dray and Manogue \cite{17}:

\begin{longtable}[]{@{}lll@{}}
\toprule\noalign{}
Algebra & Dimension & Connection to O \\
\midrule\noalign{}
\endhead
\bottomrule\noalign{}
\endlastfoot
\(G_{2}\) & 14 & Aut(O) \\
\(F_{4}\) & 52 & Aut(\(J_{3}\)(O)) \\
\(E_{6}\) & 78 & Collineations of \(OP^{2}\) \\
\(E_{8}\) & 248 & Contains all lower exceptionals \\
\end{longtable}


This chain is not accidental. It reflects the unique algebraic structure
of the octonions: Im(O) has dimension 7, the Fano plane encodes the
multiplication table, and \(G_{2}\) preserves this structure. A
\(G_{2}\)-holonomy manifold is therefore the natural geometric home for
octonionic physics, just as U(1) holonomy is the natural setting for
complex geometry.

The \(G_{2}\) structure is concretely encoded in the \textbf{standard
3-form \(\varphi\)\textsubscript{0}} on
\(\mathbb{R}\)\textsuperscript{7} (Bryant-Joyce convention):
\[\varphi_0 = e_{012} + e_{034} + e_{056} + e_{135} - e_{146} - e_{236} - e_{245}\]
where only 7 of the \(\binom{7}{3} = 35\) ordered triples carry nonzero
coefficient (all \(\pm 1\)). \(G_{2}\) is precisely the stabilizer of
\(\varphi\)\textsubscript{0} in GL(7, \(\mathbb{R}\)), and its Lie
algebra \(g_{2}\) is the kernel of the linear map
\(L_{\varphi_0} : \mathfrak{gl}(7) \to \wedge^3(\mathbb{R}^7)^*\),
\(X \mapsto \mathcal{L}_X \varphi_0\), giving
\(\dim(g_2) = 49 - \operatorname{rank}(L_{\varphi_0}) = 49 - 35 = 14\).
This is fully formalized in Lean (\texttt{G2ThreeForm.lean}): all 7
nonzero coefficients are certified by \texttt{native\_decide}, and the
Lie algebra structure (closure under addition and scalar multiplication)
is proven.

\subsection{\texorpdfstring{\(E_{8}\) x \(E_{8}\)
Structure}{E\_\{8\} x E\_\{8\} Structure}}\label{e_8-x-e_8-structure}

\(E_{8}\) is the largest exceptional simple Lie group with dimension 248
and rank 8 \cite{18}. The product \(E_{8}\) x \(E_{8}\) arises in
heterotic string theory for anomaly cancellation \cite{19}, with total
dimension 496.

The first \(E_{8}\) contains the Standard Model gauge group through the
breaking chain:

\[
E_8 \;\longrightarrow\; E_6 \times SU(3) \;\longrightarrow\; SO(10) \times U(1) \;\longrightarrow\; SU(5) \;\longrightarrow\; SU(3) \times SU(2) \times U(1)
\]

The second \(E_{8}\) provides a hidden sector whose physical
interpretation remains an open question.

Wilson (2024) demonstrates that \(E_{8}\)(-248) encodes three fermion
generations (128 degrees of freedom) with GUT structure \cite{9}. The
product dimension 496 enters the hierarchy parameter tau = (496 x
21)/(27 x 99) = 3472/891, connecting gauge structure to internal
topology.

\subsection{\texorpdfstring{The \(K_{7}\)
Manifold}{The K\_\{7\} Manifold}}\label{the-k_7-manifold}

We consider a compact, simply connected 7-manifold \(K_{7}\) with
\(G_{2}\) holonomy and Betti numbers (\(b_{2}\), \(b_{3}\)) = (21, 77).
\(G_{2}\) holonomy preserves N = 1 SUSY in 4D \cite{11}, implies Ric(g) =
0, and is the unique holonomy of \(G_{2}\) = Aut(\(\mathbb{O}\)) acting
on \(\mathbb{R}\)\textsuperscript{7}.

\textbf{Metric construction} ([A]): An explicit 7D
\(G_{2}\) metric g(s, \(\theta\), \(\psi\), y) = g\_seam(s) \(\oplus\)
g\_\{\(T^{2}\)\} \(\oplus\) g\_\{K3\}(y) on \(K_{7}\) is constructed via
Chebyshev-Cholesky parametrization (169 parameters), with K3 fiber
contribution certified at 0.07\% of total torsion [C]. The
Newton-Kantorovich certificate (h = 8.95 \(\times\) 10\(^{-9}\), margin
\(\times 56\) million, zero finite differences) establishes existence
and local uniqueness (within the NK certificate ball) of a nearby
torsion-free metric with \(\delta\)g/g \(\leq\) 4.86 \(\times\)
10\(^{-6}\). The spectral analysis ([B]) independently
confirms 21 near-zero eigenvalues of \(\Delta\)\textsubscript{2} (gap
ratio 14,635) and 77 near-zero eigenvalues of
\(\Delta\)\textsubscript{3}, consistent with the Betti numbers
(\(b_{2}\), \(b_{3}\)) = (21, 77).

\textbf{Metric decomposition}: A Chebyshev mode analysis reveals that
the metric is dominated by its constant mode (k=0), which carries
99.9998\% of the \(L^{2}\) energy. This mode is a product metric K3
\(\times\) \(T^{2}\) \(\times\) I with structural parameters g\_ss =
19/6, g\_\{\(T^{2}\)\} = 7/6, det(g) = 65/32, fixed by the structural
normalization and topological input data, with det(g) = 65/32 treated as
a metric normalization target as discussed in S1 §10.3. The k\(\geq 1\)
modes (0.0002\% of energy) constitute the minimal perturbation that
breaks the product structure and lifts the holonomy from
SU(2)\(\times\)U(1) to full \(G_{2}\). These corrections are responsible
for \(b_{1}\) = 0, the torsion ‖T‖ = 2.949\(\times 10^{-5}\), and the NK
certification, but contribute nothing to the numerical values of the
95 observables, which depend only on topological integers.

\textbf{Topological classification}: The pair (\(b_{2}\), \(b_{3}\)) =
(21, 77) does not appear among catalogued compact \(G_{2}\)
manifolds, neither in Joyce's 252 orbifold types nor in the
\textasciitilde100 CHNP TCS examples. Orthogonal TCS is excluded by
parity (\(b_{2}\)+\(b_{3}\)=98 is even; CHNP Lemma 6.7). Non-orthogonal
TCS, extra-twisted connected sums \cite{20}, and new orbifold groups
remain open paths. The NK-certified metric provides computational
evidence; a complete geometric construction is an open problem. The pair
is \textbf{unique} among all 65 known TCS literature examples (Kovalev,
CHNP, Joyce, CGN, Nordström); nearest neighbor is at distance 7.6 in
(\(b_{2}\), \(b_{3}\)) space. The certified metric is 99.98\%
block-diagonal K3\(\times T^{2} \times\)I with exact
U(1)\textsuperscript{2} symmetry (degenerate to 2\(\times 10^{-5}\)).

\textbf{Further structure}: The certified metric is smooth (no
singularities). The SM gauge group SU(3)\(\times\)SU(2)\(\times\)U(1)
does \textbf{not} arise from ADE singularities in this framework; it
emerges from the algebraic/spectral structure: \(g_{2}\) \(\subset\)
so(7) decomposition, so(8) = \(g_{2}\) \(\oplus\) L \(\oplus\) R
triality giving N\_gen = 3, and spectral gap
\(\lambda\)\textsubscript{1} = 6\(\pi\)\textsuperscript{2}/475 = 0.12467
(Richardson extrapolation: 0.12461 \(\pm\) 0.00016, deviation 0.05\%)
(§5). ADE singularities would be relevant at singular limits of
\(G_{2}\) moduli space; the smooth certified metric is at a generic
point.

The framework imposes det(g) = 65/32 = (Weyl \(\times\)
\(\alpha\)\_sum)/ 2\(^{5}\), equivalently 2 + 1/32. Joyce's theorem
\cite{20} guarantees existence of torsion-free \(G_{2}\) metrics on
suitable compact 7-manifolds; the Chebyshev-Cholesky construction
(companion paper [A]) achieves ‖T‖ \textless{} 3 \(\times\)
10\(^{-5}\) with Newton-Kantorovich certification h = 8.95 \(\times\)
10\(^{-9}\) (Chebyshev-certified, margin \(\times 56\) million; §8.5).

\subsection{Topological Constraints on Field Content}

\subsubsection{Betti Numbers as Capacity Bounds}

The Betti numbers provide upper bounds on field multiplicities:

\begin{itemize}

\item
  \textbf{\(b_{2}\)(\(K_{7}\)) = 21}: Bounds the number of gauge field
  degrees of freedom
\item
  \textbf{\(b_{3}\)(\(K_{7}\)) = 77}: Bounds the number of matter field
  degrees of freedom
\end{itemize}

\textbf{Note on gauge group origin}: In M-theory on smooth \(G_{2}\)
manifolds, dimensional reduction yields \(b_{2}\) abelian U(1) vector
multiplets \cite{11}. The non-abelian SM gauge group in GIFT emerges
instead from the algebraic/spectral structure of the \(G_{2}\) holonomy:
the \(g_{2}\) \(\subset\) so(7) decomposition and so(8) = \(g_{2}\)
\(\oplus\) L \(\oplus\) R triality (§5). The smooth certified metric is
at a generic point in \(G_{2}\) moduli space; codimension-4 ADE
singularities would be a different (singular) limit.

\subsubsection{Generation Number}

The number of chiral fermion generations follows from a topological
constraint:

\[({\rm rank}(E_8) + N_{\rm gen}) \times b_2 = N_{\rm gen} \times b_3\]

Solving: (8 + N\_gen) x 21 = N\_gen x 77 yields \textbf{N\_gen = 3}.

This derivation is formal; physically, it reflects index-theoretic
constraints on chiral zero modes, which in M-theory on \(G_{2}\) require
singular geometries for chirality \cite{25}.

\noindent

\rule{\textwidth}{0.2pt}

\section{Methodology and Epistemic Status}

\subsection{The Derivation Principle and Type Classification}

The GIFT framework derives physical observables through algebraic
combinations of 20 topological invariants (Appendix A). The 95
observables are organized by derivation directness:

\begin{longtable}[]{@{}
  >{\raggedright\arraybackslash}p{(\linewidth - 6\tabcolsep) * \real{0.1818}}
  >{\raggedright\arraybackslash}p{(\linewidth - 6\tabcolsep) * \real{0.2121}}
  >{\raggedright\arraybackslash}p{(\linewidth - 6\tabcolsep) * \real{0.3333}}
  >{\raggedright\arraybackslash}p{(\linewidth - 6\tabcolsep) * \real{0.2727}}@{}}
\toprule\noalign{}
\begin{minipage}[b]{\linewidth}\raggedright
Type
\end{minipage} & \begin{minipage}[b]{\linewidth}\raggedright
Count
\end{minipage} & \begin{minipage}[b]{\linewidth}\raggedright
Derivation
\end{minipage} & \begin{minipage}[b]{\linewidth}\raggedright
Example
\end{minipage} \\
\midrule\noalign{}
\endhead
\bottomrule\noalign{}
\endlastfoot
\textbf{I} & 33 & Direct algebra from topology & \(sin^{2} \theta\)\_W =
3/13 \\
\textbf{II} & 19 & One physical identification step & m\_u from ratio
\(\times\) VEV \\
\textbf{III} & 21 & Multi-step dynamical chains & Combined
wilson\_line+instanton lepton ratios \\
\textbf{IV} & 22 & Structural diagnostics & NK certification, Gram
conditioning \\
\end{longtable}


Type I predictions are dimensionless ratios of topological integers,
they cannot be ``fitted'' and are either correct or wrong. Type II adds
one scale identification. Type III involves dynamical mechanisms (gauge
running, eigenvalue splitting, instanton volumes). Type IV provides
internal consistency checks.

\textbf{Geometric unit convention.} Throughout this paper, we adopt the
convention that the geometric (GIFT) values are the \emph{reference},
and experimental measurements are expressed as offsets from the
geometric prediction. This reversal parallels the 2019 SI redefinition:
the kilogram is now defined from Planck's constant (exact), and any
physical realization has uncertainty relative to it. In GIFT, the
topological formula is exact; the experimental value approximates it.

\subsection{What GIFT Claims and Does Not Claim}

\textbf{Inputs}: Existence of \(K_{7}\) with \(G_{2}\) holonomy and
(\(b_{2}\), \(b_{3}\)) = (21, 77); \(E_{8} \times E_{8}\) gauge
structure; det(g) = 65/32.

\textbf{Outputs}: 95 observables across 4 types (33 Type I, 19 Type II,
21 Type III, 22 Type IV), 66 experimentally testable.

\textbf{Structure of predictions}: All 95 observables are algebraic
functions of 6 primitive topological integers (\(b_{2}\), \(b_{3}\),
dim(\(G_{2}\)), dim(\(E_{8}\)), rank(\(E_{8}\)), dim(\(K_{7}\))) plus
standard transcendentals (\(\pi\), $\sqrt{2}$, ln 2, \(\zeta\), and the golden ratio \(\varphi = (1+\sqrt{5})/2\), see §4.2.3 for the McKay $E_8 \leftrightarrow 2I$ origin of \(\varphi\)). No observable depends on
the detailed \(G_{2}\) geometry: the metric certifies that the
manifold exists but does not enter the prediction formulas. (Caveat: 6
observables use det\_num/det\_den = 65/32, which is a metric
normalization target with suggestive but not derivational algebraic
expressions in terms of topological integers, see §10.3 of S1.)

We claim that given the inputs, the outputs follow algebraically (Type
I) or computationally (Types II--IV). We do \textbf{not} claim
uniqueness of the geometry, uniqueness of the formula assignments, or
that the formula selection principle is understood.

\subsection{Three Factors Distinguishing GIFT from Numerology}

\textbf{Multiplicity}: 95 observables (66 testable), not cherry-picked
coincidences. The sensitivity analysis (§7) gives P(coincidence) =
\(10^{-346}\) under the declared uniform null model, with overdetermination ratio 2.13\(\times\), both computed
on the 33 Type I observables with experimental comparison.

\textbf{Exactness}: Several predictions are exactly rational,
\(sin^{2} \theta\)\_W = 3/13, Q\_Koide = 2/3, m\_s/m\_d = 20,
\(\Omega\)\_DM/\(\Omega\)\_b = 43/8. These cannot be fitted; they are
correct or wrong.

\textbf{Falsifiability}: DUNE will test \(\delta\)\_CP =
197\({}^{\circ}\) with expected resolution ranging from a few degrees to
\textasciitilde15\({}^{\circ}\) depending on exposure and true parameter
values (2028--2040). The SUSY spectrum (m\_gravitino = 166 GeV,
m\_moduli = 3.2 TeV) is model-dependent and constrained in standard
simplified searches; viable only in compressed or suppressed-coupling
realizations requiring dedicated recasting at HL-LHC/FCC. The proton
lifetime \(\tau\)\_p = 4.06 \(\times\) 10\(^{38}\) years exceeds
near-term Hyper-K sensitivity; Hyper-K can strengthen lower bounds and
constrain nearby GUT-scale alternatives.

\subsection{Why These Formulas?}

The selection question has two levels of answer.

\textbf{Mathematical level (deterministic).} The NK-certified metric has
zero free parameters; its 169 Chebyshev coefficients appear constrained
to a lattice generated by \{\(\pi\)\textsuperscript{2}, \(\pi\), 1, e,
\(\chi\)\textsubscript{9}\textsubscript{8}\} / (\(b_{2} \cdot b_{3}\))
(a numerical observation, not a derived result; see S1 §10.3).
Observables are functionals of this metric and are therefore algebraic in the topological invariants. Within the declared formula grammar and structural-constant set, the observed relations are highly constrained and statistically non-generic; a first-principles derivation of the full formula-selection rule remains open (see §9.4 for the extended grammar analysis with TCS atoms).

\textbf{Statistical level (empirical).} Even granting an adversary
access to the same 20 structural constants, 4.2M random algebraic
formula trees cannot reproduce GIFT's precision profile: joint
coincidence probability P = \(10^{-133}\) (algebraic null model, §7).
Structural redundancy adds a third line of evidence:
\(sin^{2} \theta\)\_W = 3/13 admits 14 independent derivations, Q\_Koide
= 2/3 admits 20 (see S2): an overdetermined web inconsistent with
post-hoc cherry-picking.

The deeper question (\emph{why \(G_{2}\) holonomy?}) has the same
epistemic status as ``why Lorentz invariance?'': \(G_{2}\) is the unique
compact exceptional holonomy group admitting a 7-dimensional
representation with the stabilizer chain \(G_{2}\) \(\supset\) SU(3)
\(\supset\) SU(2) \(\supset\) U(1) required by the Standard Model. It is
an isolated point in the space of compatible structures, not one choice
among a continuum.

\subsection{Posture: orientation, not ontology}

It is useful to distinguish three registers at which a framework of this
type operates, because the support GIFT offers is uneven across them.

\textbf{Predictive register.} GIFT specifies a finite list of inputs and
derives 95 observables from them; 66 of these are testable, with mean
deviation \textasciitilde1\% and a sensitivity profile that places the
joint configuration at \(> 4.5\sigma\) against random algebraic null
models (§7). This is the register where the framework either holds or
fails, and it is the only register at which we make any positive
claim.

\textbf{Architectural register.} The choice of \(G_{2}\)-holonomy
7-manifolds with \(E_8 \times E_8\) structure is \emph{motivated} by
mathematical considerations (§§2, 9.4): \(G_{2}\) is the unique
exceptional compact holonomy admitting the SU(3) \(\times\) SU(2)
\(\times\) U(1) stabilizer chain, \(E_8\) is the largest exceptional Lie
algebra, and the topological constraints fix \((b_2, b_3) = (21, 77)\)
and \(N_{\mathrm{gen}} = 3\). We argue this architecture is
\emph{natural}; we do not argue it is \emph{necessary}.

\textbf{Ontological register.} A reader may wonder whether GIFT carries
a thesis about reality: for instance, the speculation that geometry,
information, and energy are not three correlated aspects of nature but
three views of a single underlying configuration, or the resonance with
Wheeler's ``It from Bit'' programme \cite{39} and the holographic
principle. Such readings are compatible with the framework and
historically motivated its development, but they are \textbf{neither
required nor demonstrated} by the predictions. A reader who finds the
Wheeler-holographic picture persuasive will see GIFT as a natural piece
of that puzzle; a reader who prefers a strictly empirical stance will
see a falsifiable predictive framework. Both readings are defensible;
the framework requires neither.

In Worrall's \cite{40} and Ladyman's \cite{41} terms, GIFT is best read as
a \emph{moderate structural-realist orientation}: structure carries
predictive weight independently of any further ontological commitment.
The framework's success or failure is decided by experiment in the
predictive register, not by adjudication in the ontological one. A more
extended discussion of this posture, written for a non-technical
audience, appears in the companion essay \emph{Orientation, not
ontology} [essay].

\noindent

\rule{\textwidth}{0.2pt}

\section{Observables: 95 Relations from 20 Structural Constants}

\subsection{Gauge Sector}

\subsubsection{Weinberg Angle}

\[\sin^2\theta_W = \frac{b_2}{b_3 + \dim(G_2)} = \frac{21}{91} = \frac{3}{13} = 0.230769\]

Experimental (PDG 2024) \cite{1}: 0.23122 +/- 0.00004. Deviation:
\textbf{0.195\%}.

The numerator \(b_{2}\) counts gauge moduli; the denominator \(b_{3}\) +
dim(\(G_{2}\)) counts matter plus holonomy degrees of freedom. The ratio
measures gauge-matter coupling geometrically.

\subsubsection{Strong Coupling}

\[\alpha_s(M_Z) = \frac{\sqrt{2}}{\dim(G_2) - p_2} = \frac{\sqrt{2}}{12} = 0.11785\]

Experimental: 0.1180 +/- 0.0009. Deviation: \textbf{0.126\%}.

\subsection{Lepton Sector}

\subsubsection{Koide Parameter}

The Koide formula has resisted explanation since 1982. Koide discovered
an empirical relation among the charged lepton masses \cite{6}:

\[Q = \frac{(m_e + m_\mu + m_\tau)^2}{(\sqrt{m_e} + \sqrt{m_\mu} + \sqrt{m_\tau})^2} = \frac{2}{3}\]

Using contemporary mass values, this holds to six significant figures:
Q\_exp = 0.666661 +/- 0.000007.

GIFT provides:

\[Q_{\rm Koide} = \frac{\dim(G_2)}{b_2} = \frac{14}{21} = \frac{2}{3}\]

The derivation requires only two topological invariants: dim(\(G_{2}\))
= 14 (holonomy group dimension) and \(b_{2}\) = 21 (second Betti
number). No fitting is involved.

\begin{longtable}[]{@{}
  >{\raggedright\arraybackslash}p{(\linewidth - 4\tabcolsep) * \real{0.3846}}
  >{\raggedright\arraybackslash}p{(\linewidth - 4\tabcolsep) * \real{0.3077}}
  >{\raggedright\arraybackslash}p{(\linewidth - 4\tabcolsep) * \real{0.3077}}@{}}
\toprule\noalign{}
\begin{minipage}[b]{\linewidth}\raggedright
Approach
\end{minipage} & \begin{minipage}[b]{\linewidth}\raggedright
Result
\end{minipage} & \begin{minipage}[b]{\linewidth}\raggedright
Status
\end{minipage} \\
\midrule\noalign{}
\endhead
\bottomrule\noalign{}
\endlastfoot
Preon models (Koide 1982) & Q = 2/3 assumed & Circular \\
\(S_{3}\) symmetry (various) & Q \textasciitilde{} 2/3 fitted &
Approximate \\
\textbf{GIFT} & \textbf{Q = dim(\(G_{2}\))/\(b_{2}\) = 14/21 = 2/3} &
\textbf{Algebraic identity} \\
\end{longtable}


Deviation: \textbf{0.0009\%}. This is the most precise agreement in the
framework.

\subsubsection{Tau-Electron Mass Ratio}

\[\frac{m_\tau}{m_e} = \dim(K_7) + 10 \times \dim(E_8) + 10 \times H^* = 7 + 2480 + 990 = 3477\]

Experimental: 3477.15 +/- 0.05. Deviation: \textbf{0.004\%}.

The integer 3477 = 3 x 19 x 61 = N\_gen x prime(8) x \$kappa\_T\^{}-\$1
factorizes into framework constants.

\subsubsection{Muon-Electron Mass Ratio}

\[\frac{m_\mu}{m_e} = \dim(J_3(\mathbb{O}))^\phi = 27^\phi = 207.01\]

where \(\varphi = (1+\sqrt{5})/2\) is the golden ratio, arising from the McKay
correspondence \(E_{8}\) \(\leftrightarrow\) 2I (binary icosahedral
group), which links \(E_{8}\) to icosahedral geometry. Experimental:
206.768. Deviation: \textbf{0.118\%}.

\textbf{Note}: \(\varphi\) is the only non-integer input among the 33
Type I predictions and does not appear in the 20 structural constants of
§S3.3. Its status is accordingly weaker than the other Type I
derivations: the formula is algebraically exact and Lean-certified, but
\(\varphi\)'s derivation from \(E_{8}\) structure requires the
additional McKay step (see S2 §11).

\subsection{Quark Sector}

\[\frac{m_s}{m_d} = p_2^2 \times \text{Weyl} = 4 \times 5 = 20\]

Experimental (PDG 2024): 20.0 +/- 1.0. Deviation: \textbf{0.00\%}.

\[\frac{m_b}{m_t} = \frac{b_0}{2b_2} = \frac{1}{42}\]

The constant 42 = \(p_{2}\) x N\_gen x dim(\(K_{7}\)) = 2 x 3 x 7 is a
structural invariant (not to be confused with chi(\(K_{7}\)) = 0, which
vanishes for any compact odd-dimensional manifold).

Experimental: 0.024 +/- 0.001. Deviation: \textbf{0.79\%}.

\subsection{Neutrino Sector}

\subsubsection{CP-Violation Phase}

The GIFT prediction for the CP-violation phase is:

\[\delta_{CP} = \dim(K_7) \times \dim(G_2) + H^* = 7 \times 14 + 99 = 197°\]

decomposing into a local contribution (7 \(\times\) 14 = 98,
fiber-holonomy coupling) and a global contribution (H* = 99,
cohomological dimension). This is the canonical GIFT prediction, pure
topological, zero corrections.

\textbf{Experimental status}: NuFIT 5.2 (2022) gave \(\delta\)\_CP
\(\approx\) 197\({}^{\circ}\), an exact match. NuFIT 6.0 (Oct 2024)
shifted the central value to 177\({}^{\circ}\) \(\pm\) 20\({}^{\circ}\),
creating an 11.3\% deviation. The experimental uncertainty is large
(\(\pm 20\)\({}^{\circ}\)), and the central value may shift further as
T2K, NOvA, and DUNE accumulate statistics.

A post-hoc structural observation involving a possible compactification factor 62/69 is recorded in S2 Appendix F for completeness; it is \textbf{not} part of the main framework and is not used to revise the canonical 197\({}^{\circ}\) prediction.

\textbf{Falsification criterion}: If DUNE measures \(\delta\)\_CP
outside {[}182, 212{]}\({}^{\circ}\) at 3\(\sigma\), the framework faces
serious tension. The 197\({}^{\circ}\) prediction sits at the edge of
the current NuFIT 6.0 1\(\sigma\) band (177\({}^{\circ}\) \(\pm\)
20\({}^{\circ}\) = {[}157\({}^{\circ}\), 197\({}^{\circ}\){]}).

\subsubsection{Mixing Angles}

\begin{longtable}[]{@{}
  >{\raggedright\arraybackslash}p{(\linewidth - 8\tabcolsep) * \real{0.1591}}
  >{\raggedright\arraybackslash}p{(\linewidth - 8\tabcolsep) * \real{0.2045}}
  >{\raggedright\arraybackslash}p{(\linewidth - 8\tabcolsep) * \real{0.1364}}
  >{\raggedright\arraybackslash}p{(\linewidth - 8\tabcolsep) * \real{0.3636}}
  >{\raggedright\arraybackslash}p{(\linewidth - 8\tabcolsep) * \real{0.1364}}@{}}
\toprule\noalign{}
\begin{minipage}[b]{\linewidth}\raggedright
Angle
\end{minipage} & \begin{minipage}[b]{\linewidth}\raggedright
Formula
\end{minipage} & \begin{minipage}[b]{\linewidth}\raggedright
GIFT
\end{minipage} & \begin{minipage}[b]{\linewidth}\raggedright
NuFIT 6.0 \cite{27}
\end{minipage} & \begin{minipage}[b]{\linewidth}\raggedright
Dev.
\end{minipage} \\
\midrule\noalign{}
\endhead
\bottomrule\noalign{}
\endlastfoot
theta\_12 & arctan(sqrt(delta/gamma\_GIFT)) & 33.40 deg & 33.41 +/- 0.75
deg & 0.03\% \\
theta\_13 & pi/\(b_{2}\) & 8.57 deg & 8.54 +/- 0.12 deg & 0.37\% \\
theta\_23 & arcsin((\(b_{3}\) - \(p_{2}\))/H*) & 49.25 deg & 49.3 +/-
1.0 deg & 0.10\% \\
\end{longtable}


The auxiliary parameters: delta = 2\emph{pi/\(Weyl^{2}\) = 2}pi/25 and
gamma\_GIFT = (2 x rank(\(E_{8}\)) + 5 x H*)/(10 x dim(\(G_{2}\)) + 3 x
dim(\(E_{8}\))) = 511/884.

\subsection{Higgs Sector}

\[\lambda_H = \frac{\sqrt{\dim(G_2) + N_{gen}}}{\det(g)_{den}} = \frac{\sqrt{17}}{32} = 0.1289\]

The Higgs self-coupling combines holonomy dimension with generation
count, normalized by the metric determinant scale. Experimental (PDG
2024): 0.129 \(\pm\) 0.001. Deviation: \textbf{0.12\%}. A TCS
alternative \(\lambda\)\_H =
\(b_{2}\)(\(M_{1}\))/(\(b_{3}\)+\(b_{2}\)(\(M_{2}\))) = 11/87 = 0.1264
is purely rational (see S2 §17).

\subsection{Boson Mass Ratios}

\begin{longtable}[]{@{}
  >{\raggedright\arraybackslash}p{(\linewidth - 8\tabcolsep) * \real{0.2553}}
  >{\raggedright\arraybackslash}p{(\linewidth - 8\tabcolsep) * \real{0.1915}}
  >{\raggedright\arraybackslash}p{(\linewidth - 8\tabcolsep) * \real{0.1277}}
  >{\raggedright\arraybackslash}p{(\linewidth - 8\tabcolsep) * \real{0.2979}}
  >{\raggedright\arraybackslash}p{(\linewidth - 8\tabcolsep) * \real{0.1277}}@{}}
\toprule\noalign{}
\begin{minipage}[b]{\linewidth}\raggedright
Observable
\end{minipage} & \begin{minipage}[b]{\linewidth}\raggedright
Formula
\end{minipage} & \begin{minipage}[b]{\linewidth}\raggedright
GIFT
\end{minipage} & \begin{minipage}[b]{\linewidth}\raggedright
Experimental
\end{minipage} & \begin{minipage}[b]{\linewidth}\raggedright
Dev.
\end{minipage} \\
\midrule\noalign{}
\endhead
\bottomrule\noalign{}
\endlastfoot
m\_H/m\_W & (N\_gen + dim(\(E_{6}\)))/dim(\(F_{4}\)) = 81/52 & 1.5577 &
1.558 +/- 0.002 & 0.02\% \\
m\_W/m\_Z & (2\(b_{2}\) - Weyl)/(2\(b_{2}\)) = 37/42 & 0.8810 & 0.8815
+/- 0.0002 & 0.06\% \\
m\_H/m\_t & fund(\(E_{7}\))/\(b_{3}\) = 56/77 & 0.7273 & 0.725 +/- 0.003
& 0.31\% \\
\end{longtable}


\subsection{CKM Matrix}

\begin{longtable}[]{@{}
  >{\raggedright\arraybackslash}p{(\linewidth - 8\tabcolsep) * \real{0.2553}}
  >{\raggedright\arraybackslash}p{(\linewidth - 8\tabcolsep) * \real{0.1915}}
  >{\raggedright\arraybackslash}p{(\linewidth - 8\tabcolsep) * \real{0.1277}}
  >{\raggedright\arraybackslash}p{(\linewidth - 8\tabcolsep) * \real{0.2979}}
  >{\raggedright\arraybackslash}p{(\linewidth - 8\tabcolsep) * \real{0.1277}}@{}}
\toprule\noalign{}
\begin{minipage}[b]{\linewidth}\raggedright
Observable
\end{minipage} & \begin{minipage}[b]{\linewidth}\raggedright
Formula
\end{minipage} & \begin{minipage}[b]{\linewidth}\raggedright
GIFT
\end{minipage} & \begin{minipage}[b]{\linewidth}\raggedright
Experimental
\end{minipage} & \begin{minipage}[b]{\linewidth}\raggedright
Dev.
\end{minipage} \\
\midrule\noalign{}
\endhead
\bottomrule\noalign{}
\endlastfoot
\(sin^{2}\)(theta\_12\_CKM) & fund(\(E_{7}\))/dim(\(E_{8}\)) = 56/248 &
0.2258 & 0.2250 +/- 0.0006 & 0.36\% \\
A\_Wolfenstein & (Weyl + dim(\(E_{6}\)))/H* = 83/99 & 0.838 & 0.836 +/-
0.015 & 0.29\% \\
\(sin^{2}\)(theta\_23\_CKM) & dim(\(K_{7}\))/PSL(2,7) = 7/168 & 0.0417 &
0.0412 +/- 0.0008 & 1.13\% \\
\end{longtable}


The Cabibbo angle emerges from the ratio of the \(E_{7}\) fundamental
representation to \(E_{8}\) dimension.

\subsection{Cosmological Observables}

\begin{longtable}[]{@{}
  >{\raggedright\arraybackslash}p{(\linewidth - 8\tabcolsep) * \real{0.2553}}
  >{\raggedright\arraybackslash}p{(\linewidth - 8\tabcolsep) * \real{0.1915}}
  >{\raggedright\arraybackslash}p{(\linewidth - 8\tabcolsep) * \real{0.1277}}
  >{\raggedright\arraybackslash}p{(\linewidth - 8\tabcolsep) * \real{0.2979}}
  >{\raggedright\arraybackslash}p{(\linewidth - 8\tabcolsep) * \real{0.1277}}@{}}
\toprule\noalign{}
\begin{minipage}[b]{\linewidth}\raggedright
Observable
\end{minipage} & \begin{minipage}[b]{\linewidth}\raggedright
Formula
\end{minipage} & \begin{minipage}[b]{\linewidth}\raggedright
GIFT
\end{minipage} & \begin{minipage}[b]{\linewidth}\raggedright
Experimental
\end{minipage} & \begin{minipage}[b]{\linewidth}\raggedright
Dev.
\end{minipage} \\
\midrule\noalign{}
\endhead
\bottomrule\noalign{}
\endlastfoot
Omega\_DM/Omega\_b & (1 + 2\(b_{2}\))/rank(\(E_{8}\)) = 43/8 & 5.375 &
5.375 +/- 0.1 & 0.00\% \\
n\_s & zeta(11)/zeta(5) & 0.9649 & 0.9649 +/- 0.0042 & 0.004\% \\
h (Hubble) & (PSL(2,7) - 1)/dim(\(E_{8}\)) = 167/248 & 0.6734 & 0.674
+/- 0.005 & 0.09\% \\
Omega\_b/Omega\_m & Weyl/det(g)\_den = 5/32 & 0.1562 & 0.157 +/- 0.003 &
0.16\% \\
sigma\_8 & (\(p_{2}\) + 32)/(2\(b_{2}\)) = 34/42 & 0.8095 & 0.811 +/-
0.006 & 0.18\% \\
Omega\_DE & ln(2)\(\times\)(\(b_{2}\)+\(b_{3}\))/H* =
ln(2)\(\times 98\)/99 & 0.6861 & 0.6847 +/- 0.005 & 0.21\% \\
Y\_p & (1 + dim(\(G_{2}\)))/\$kappa\_T\^{}-\$1 = 15/61 & 0.2459 & 0.245
+/- 0.003 & 0.37\% \\
\end{longtable}


The dark-to-baryonic matter ratio Omega\_DM/Omega\_b = 43/8 is exact.
The structural invariant 2\(b_{2}\) = 42 that gives m\_b/m\_t = 1/42
also determines this cosmological ratio, connecting quark physics to
large-scale structure through \(K_{7}\) geometry.

\subsection{Type I Summary (33 Observables)}

The 33 Type I predictions derive directly from the 20 structural
constants (Appendix A). All are dimensionless ratios; all 33 are
Lean-certified. Representative highlights:

\begin{longtable}[]{@{}
  >{\raggedright\arraybackslash}p{(\linewidth - 10\tabcolsep) * \real{0.20}}
  >{\raggedright\arraybackslash}p{(\linewidth - 10\tabcolsep) * \real{0.22}}
  >{\raggedright\arraybackslash}p{(\linewidth - 10\tabcolsep) * \real{0.10}}
  >{\raggedright\arraybackslash}p{(\linewidth - 10\tabcolsep) * \real{0.12}}
  >{\raggedright\arraybackslash}p{(\linewidth - 10\tabcolsep) * \real{0.10}}
  >{\raggedright\arraybackslash}p{(\linewidth - 10\tabcolsep) * \real{0.26}}@{}}
\toprule\noalign{}
\begin{minipage}[b]{\linewidth}\raggedright
Observable
\end{minipage} & \begin{minipage}[b]{\linewidth}\raggedright
Formula
\end{minipage} & \begin{minipage}[b]{\linewidth}\raggedright
GIFT
\end{minipage} & \begin{minipage}[b]{\linewidth}\raggedright
Exp.
\end{minipage} & \begin{minipage}[b]{\linewidth}\raggedright
Dev.
\end{minipage} & \begin{minipage}[b]{\linewidth}\raggedright
Status
\end{minipage} \\
\midrule\noalign{}
\endhead
\bottomrule\noalign{}
\endlastfoot
\(sin^{2} \theta\)\_W & \(b_{2}\)/(\(b_{3}\)+dim(\(G_{2}\))) = 21/91 &
3/13 & 0.23122 & 0.195\% & TOPOLOGICAL \\
Q\_Koide & dim(\(G_{2}\))/\(b_{2}\) = 14/21 & 2/3 & 0.666661 &
0.001\% & TOPOLOGICAL \\
\(\alpha\)\textsuperscript{-}\textsuperscript{1} & 128 + 9 +
det(g)\(\times \kappa\)\_T & 137.033 & 137.036 & 0.002\% & TOPOLOGICAL (uses metric normalization) \\
m\_\(\tau\)/m\_e & 7 + 2480 + 990 & 3477 & 3477.15 & 0.004\% & TOPOLOGICAL \\
n\_s & \(\zeta\)(11)/\(\zeta\)(5) & 0.9649 & 0.9649 & 0.004\% & DERIVED \\
\(\Omega\)\_DM/\(\Omega\)\_b & (1+42)/8 & 43/8 & 5.375 & 0.00\% & TOPOLOGICAL \\
\end{longtable}

All 33 Type I observables are VERIFIED in Lean 4. Status legend: \textbf{TOPOLOGICAL} = pure topological-integer ratio; \textbf{DERIVED} = closed-form involving standard transcendentals (\(\pi\), \(\sqrt{2}\), \(\ln 2\), \(\zeta\), golden ratio \(\varphi\)); \textbf{STRUCTURAL} = depends on the imposed metric normalization (det\_num / det\_den).

Complete derivations for all 33 in Supplement S2; full 95-entry table in
Supplement S3.

\subsection{Type II: Extended Algebraic Predictions (19)}

Type II observables require one physical identification step beyond Type
I ratios. Representative results:

\textbf{Absolute quark masses} (via m\_q = ratio \(\times\)
reference\_mass): m\_u = 2.16 MeV (0.00\%), m\_d = 4.67 MeV (0.22\%),
m\_s = 93.4 MeV (0.22\%), m\_c = 1.27 GeV (0.00\%), m\_b = 4.18 GeV
(0.00\%), m\_t = 172.7 GeV (0.01\%).

\textbf{CKM magnitudes} (from Wolfenstein parametrization):
\textbar V\_us\textbar{} = 0.2253 (0.13\%), \textbar V\_cb\textbar{} =
0.0412 (0.24\%), \textbar V\_ub\textbar{} = 0.00365 (0.27\%),
\textbar V\_td\textbar{} = 0.0087 (0.34\%), \textbar V\_tb\textbar{} =
0.9991 (0.00\%).

\textbf{Extended ratios}: m\_c/m\_s = 246/21 = 11.714 (exp 11.7, dev
0.12\%), m\_c/m\_d = 234.3 (exp 234.0, dev 0.12\%),
m\_\(\mu\)/m\_\(\tau\) = 5/84 (exp 0.0595, dev 0.04\%).

Type II mean deviation: \textbf{0.17\%} across 19 observables. These
inherit the precision of Type I ratios, with the physical identification
step contributing negligible additional error.

\subsection{Type III: Dynamical Predictions (21)}

Type III observables involve multi-step dynamical mechanisms. They are
grouped by computation:

\textbf{wilson\_line Non-adiabatic} (3 obs): Wilson line eigenvalue
splitting on K3 fiber at c = 0.452 gives raw lepton mass ratios with
0.5--2.1\% deviation. These are improved to \textless{} 0.4\% by the
combined wilson\_line+instanton pipeline (§6.4).

\textbf{RGE\_running RGE running} (4 obs): Two-loop MSSM evolution from
M\_GUT to M\_Z. The topological \(sin^{2} \theta\)\_W = 3/13 at M\_GUT
runs to 0.2377 at M\_Z (exp 0.2312, dev 2.78\%). The strong coupling
\(\alpha\)\_s(RGE) = 0.1224 deviates 3.7\% using \(G_{2}\)-MSSM
split-spectrum matching (§5.3). M\_GUT = 2 \(\times\) 10\(^{16}\) GeV is
an exact match.

\textbf{spectral Spectral} (5 obs): effective Weyl law exponent (from
the adiabatic seam-sector decomposition, extrapolated to d\_eff = 7)
\(\alpha\) = 3.460 (exp 3.5, dev 1.1\%), 22,671 KK states below cutoff,
57,578 fiber channels, Poisson level spacing. See [B] for
the direct seam-sector result \(\alpha\) = 1.998.

\textbf{gauge\_bundle Gauge bundle} (4 obs): cond(f\_IJ) = 1.047
(near-perfect gauge universality), \(\alpha\)\_ratio = 1.000002,
effective Yukawa rank = 3, and \(\kappa\)(gauge) = 1.047 (4.7\%
departure from exact universality, the largest Type III bundle
deviation).

\textbf{instanton Instanton + combined} (5 obs): Associative volume
differences give \(\Delta\)V(e-\(\tau\)) = 8.633 (dev 5.9\%),
\(\Delta\)V(e-\(\mu\)) = 3.271 (dev 15.9\%). Combined
wilson\_line+instanton pipeline with \(\alpha\) = \(e^K\) (geometric,
§6.4) gives \(\tau\)/e = 3485 (0.24\%), \(\tau\)/\(\mu\) = 16.69
(0.75\%), \(\mu\)/e = 208.8 (0.97\%).

Type III mean deviation: \textbf{3.4\%} across 14 experimentally
comparable observables (2.3\% across all 21, median 1.6\%). Details in
§5 (Gauge), §6 (Mass Hierarchy). (M\_res and N\_QNM from Pinčák et
al.~2026 \cite{42} are classified Type IV, see §4.11.)

\subsection{Type IV: Structural Diagnostics (22)}

Type IV observables are internal consistency checks with no experimental
comparison:

\begin{itemize}

\item
  \textbf{Topology}: \(b_{2}\) = 21, \(b_{3}\) = 77, \(\chi\)(\(K_{7}\))
  = 0, H* = 99
\item
  \textbf{Newton-Kantorovich certification}: h = 8.95 \(\times\)
  10\(^{-9}\) (\textless{} 0.5, margin \(\times 56\)M), \(\delta\)g/g
  \(\leq\) 4.86 \(\times\) 10\(^{-6}\)
\item
  \textbf{Gram conditioning}: cond(G\_K3) = 1.05, cond(G\_K7) = 1.05,
  cond(G\_35) = 7.66, G\_77 positive definite
\item
  \textbf{Spectral counts}: 22,671 KK states, 57,578 fiber channels,
  Poisson spacing confirmed
\item
  \textbf{Torsion}: ‖T‖\_\(C^{0}\) = 2.949 \(\times\) 10\(^{-5}\)
  (\(\times 2995\) reduction in 5 Joyce steps)
\item
  \textbf{Metric eigenvalues}: g\_ss = 19/6, g\_\{\(T^{2}\)\} = 7/6,
  g\_\{K3\} \(\approx\) 64/77
\item
  \textbf{Instanton \& BH diagnostics} (Pinčák et al.~2026 \cite{42}):
  N\_QNM = 98 QNM mode families,
  \(b_{3}\)/\(b_{3}\)(\(S^{3} \times S^{4}\)) = 77\(\times\) instanton
  suppression, M\_res = v\_\(EW^{2}\)/M\_Pl (BH remnant mass: no
  experimental comparison)
\end{itemize}

These diagnostics confirm the geometric construction is well-conditioned
and internally consistent.

\subsection{Summary Statistics (All 95 Observables)}

\textbf{Global performance} (95 observables):

\begin{longtable}[]{@{}lllll@{}}
\toprule\noalign{}
Metric & Type I & Type II & Type III & All (I+II+III) \\
\midrule\noalign{}
\endhead
\bottomrule\noalign{}
\endlastfoot
Count & 33 & 19 & 21 & 73* \\
With exp. comparison & 33 & 19 & 14 & 66 \\
Mean deviation & 0.73\% & 0.17\% & 3.44\% & \textasciitilde1.15\% \\
Median deviation & 0.23\% & 0.12\% & 1.39\% & 0.23\% \\
Exact matches (\textless0.01\%) & 5 & 3 & 3 & 11 \\
Within 1\% & 28 & 19 & 6 & 53 \\
Maximum deviation & 11.3\% & 0.79\% & 15.9\% & 15.9\% \\
\end{longtable}


\(\dagger \delta\)\_CP = 11.3\% dominates Type I; excluding it, Type I
mean = 0.40\% (28 within-1\%, 5 exact matches, max 2.77\%).

*66 with experimental comparison out of 73 (I+II+III); 22 Type IV
observables are structural.

\textbf{Sector breakdown} (11 sectors; 41 observables listed, see
§5--6 for remaining 25 comparable):

\begin{longtable}[]{@{}
  >{\raggedright\arraybackslash}p{(\linewidth - 8\tabcolsep) * \real{0.2105}}
  >{\raggedright\arraybackslash}p{(\linewidth - 8\tabcolsep) * \real{0.1842}}
  >{\raggedright\arraybackslash}p{(\linewidth - 8\tabcolsep) * \real{0.2632}}
  >{\raggedright\arraybackslash}p{(\linewidth - 8\tabcolsep) * \real{0.1579}}
  >{\raggedright\arraybackslash}p{(\linewidth - 8\tabcolsep) * \real{0.1842}}@{}}
\toprule\noalign{}
\begin{minipage}[b]{\linewidth}\raggedright
Sector
\end{minipage} & \begin{minipage}[b]{\linewidth}\raggedright
N\_obs
\end{minipage} & \begin{minipage}[b]{\linewidth}\raggedright
Mean dev
\end{minipage} & \begin{minipage}[b]{\linewidth}\raggedright
Best
\end{minipage} & \begin{minipage}[b]{\linewidth}\raggedright
Worst
\end{minipage} \\
\midrule\noalign{}
\endhead
\bottomrule\noalign{}
\endlastfoot
Electroweak & 3 & 0.11\% &
\(\alpha\)\textsuperscript{-}\textsuperscript{1} 0.002\% &
\(sin^{2} \theta\)\_W 0.20\% \\
Boson & 3 & 0.13\% & m\_H/m\_W 0.02\% & m\_H/m\_t 0.31\% \\
Lepton & 3 & 0.04\% & Q\_Koide 0.001\% & m\_\(\mu\)/m\_e 0.12\% \\
Quark & 4 & 0.21\% & m\_s/m\_d 0.00\% & m\_b/m\_t 0.79\% \\
Cosmology & 7 & 0.15\% & n\_s 0.004\% & Y\_p 0.37\% \\
PMNS & 4 & 0.29\% & \(\theta\)\textsubscript{1}\textsubscript{2} 0.03\%
& \(\theta\)\textsubscript{1}\textsubscript{3} 0.37\% \\
CKM & 3 & 0.59\% & A\_Wolf 0.29\% &
\(sin^{2} \theta\)\textsubscript{2}\textsubscript{3} 1.13\% \\
Gauge (running) & 4 & 2.3\% & M\_GUT 0.00\% & \(\alpha\)\_s(RGE)
3.7\% \\
Instanton & 3 & 7.4\% & \(\Delta\)V(e-\(\tau\)) 5.9\% &
\(\Delta\)V(e-\(\mu\)) 15.9\% \\
Combined & 3 & 0.66\% & \(\tau\)/e 0.23\% & \(\mu\)/e 0.98\% \\
Bundle & 4 & 1.6\% & \(\alpha\)\_ratio 0.00\% & \(\kappa\)(gauge)
4.7\% \\
\end{longtable}


\noindent

\rule{\textwidth}{0.2pt}

\section{\texorpdfstring{Gauge Sector: From \(E_{8}\) to the Standard
Model}{Gauge Sector: From E\_\{8\} to the Standard Model}}\label{gauge-sector-from-e_8-to-the-standard-model}

The gauge sector derives the Standard Model gauge group from the
\(E_{8} \times E_{8}\) structure of heterotic M-theory on \(K_{7}\).
This section presents the complete breaking chain, anomaly cancellation,
gauge coupling running, and bundle universality, all from the
topological data of \(K_{7}\) and an explicit 7D \(G_{2}\) metric (169
optimized Chebyshev-Cholesky parameters capturing the dominant seam
sector, K3 fiber certified at 0.07\% {[}30, 45{]}), with no free
parameters in the physical predictions.

\subsection{\texorpdfstring{The \(E_{8}\) \(\to\) Standard Model
Breaking
Chain}{The E\_\{8\} \textbackslash to Standard Model Breaking Chain}}\label{the-e_8-to-standard-model-breaking-chain}

The first \(E_{8}\) factor breaks to the Standard Model through a
six-level chain:

\begin{longtable}[]{@{}
  >{\raggedright\arraybackslash}p{(\linewidth - 8\tabcolsep) * \real{0.1628}}
  >{\raggedright\arraybackslash}p{(\linewidth - 8\tabcolsep) * \real{0.1628}}
  >{\raggedright\arraybackslash}p{(\linewidth - 8\tabcolsep) * \real{0.2558}}
  >{\raggedright\arraybackslash}p{(\linewidth - 8\tabcolsep) * \real{0.2558}}
  >{\raggedright\arraybackslash}p{(\linewidth - 8\tabcolsep) * \real{0.1628}}@{}}
\toprule\noalign{}
\begin{minipage}[b]{\linewidth}\raggedright
Level
\end{minipage} & \begin{minipage}[b]{\linewidth}\raggedright
Group
\end{minipage} & \begin{minipage}[b]{\linewidth}\raggedright
Dimension
\end{minipage} & \begin{minipage}[b]{\linewidth}\raggedright
Mechanism
\end{minipage} & \begin{minipage}[b]{\linewidth}\raggedright
Scale
\end{minipage} \\
\midrule\noalign{}
\endhead
\bottomrule\noalign{}
\endlastfoot
0 & \(E_{8} \times E_{8}\) & 496 & Heterotic structure & M\_Pl \\
1 & \(E_{8}\) & 248 & Second \(E_{8}\) = hidden sector & M\_string \\
2 & \(E_{6}\) \(\times\) SU(3) & 78 + 8 = 86 & Adjoint branching &
M\_string \\
3 & SO(10) \(\times\) U(1) & 45 + 1 = 46 & \(E_{6}\) breaking &
M\_GUT \\
4 & SU(5) \(\times\) U(1) & 24 + 1 = 25 & Pati-Salam intermediate &
M\_GUT \\
5 & SU(3) \(\times\) SU(2) \(\times\) U(1) & 8 + 3 + 1 = 12 & Standard
Model & M\_Z \\
\end{longtable}


The \(E_{8}\) adjoint decomposes under \(E_{6}\) \(\times\) SU(3) as:

\[248 = (78, 1) + (1, 8) + (27, 3) + (\overline{27}, \overline{3})\]

yielding 78 + 8 + 81 + 81 = 248. The fundamental of SU(3) has dimension
3, giving \textbf{N\_gen = 3} chiral families from the (27, 3)
representation. This reproduces the index-theorem result of §2.

\textbf{Fundamental group}: \(\pi\)\textsubscript{1}(\(K_{7}\)) = \{1\}
(simply connected, from \(b_{1}\) = 0 for \(G_{2}\) holonomy).
Traditional Wilson line breaking via \(\pi\)\textsubscript{1} is
therefore trivial. Instead, the breaking proceeds through the \(Z_{3}\)
lattice action on the K3 fiber (§6.1).

\textbf{Scales}: M\_Pl = 1.22 \(\times\) 10\(^{19}\) GeV, M\_string = 4
\(\times\) 10\(^{17}\) GeV, M\_GUT = 2 \(\times\) 10\(^{16}\) GeV, M\_Z
= 91.19 GeV. The scale hierarchy M\_GUT/M\_Z \textasciitilde{} 2
\(\times\) 10\(^{14}\) emerges from the geometry without fine-tuning.

\subsection{Anomaly Cancellation}

All six Standard Model gauge anomalies vanish:

\begin{longtable}[]{@{}lll@{}}
\toprule\noalign{}
Anomaly & Value & Status \\
\midrule\noalign{}
\endhead
\bottomrule\noalign{}
\endlastfoot
SU(3)\textsuperscript{3} & 0.0 & Exact \\
SU(2)\textsuperscript{3} & 0.0 & Exact \\
U(1)\textsuperscript{3} & 2.3 \(\times\) 10\(^{-16}\) & Machine zero \\
grav-U(1) & 0.0 & Exact \\
SU(3)\textsuperscript{2}-U(1) & 0.0 & Exact \\
SU(2)\textsuperscript{2}-U(1) & 0.0 & Exact \\
\end{longtable}


The Green-Schwarz mechanism provides additional consistency: -
\(\chi\)(\(K_{7}\)) = 0 (tadpole cancellation for compact
odd-dimensional manifolds) - Tadpole value = 0 (no net D-brane charge) -
\(G_{4}\) flux lattice rank = 77 = \(b_{3}\)(\(K_{7}\)) (all flux
degrees of freedom realized)

The 10/10 verification certificate confirms all anomaly, tadpole, and
lattice checks pass. Lean: \texttt{TCSGaugeBreaking.lean} (10 conjuncts,
0 sorry).

\subsection{Gauge Coupling Running}

\textbf{Boundary conditions at M\_GUT}: \(\alpha\)\_\(GUT^{-1}\) = 25.3,
\(sin^{2} \theta\)\_W = 3/13 = 0.23077. These are topological:
\(\alpha\)\_GUT derives from the gauge kinetic function on \(K_{7}\),
and \(sin^{2} \theta\)\_W from the Betti number ratio
\(b_{2}\)/(\(b_{3}\) + dim(\(G_{2}\))).

\textbf{Two-loop MSSM RGE} (tan \(\beta\) = 2, M\_SUSY = m\_moduli =
3165 GeV):

\begin{longtable}[]{@{}llll@{}}
\toprule\noalign{}
Coupling & GIFT at M\_Z & Experimental & Deviation \\
\midrule\noalign{}
\endhead
\bottomrule\noalign{}
\endlastfoot
\(sin^{2} \theta\)\_W & 0.2377 & 0.23122 & 2.78\% \\
\(\alpha\)\_\(em^{-1}\) & 131.19 & 127.951 & 2.53\% \\
\(\alpha\)\_s (all-MSSM) & 0.1038 & 0.1180 & 12.1\% \\
**\(\alpha\)\_s (split-spectrum)** & \textbf{0.1224} & \textbf{0.1180} &
\textbf{3.7\%} \\
\end{longtable}


A naive degenerate-spectrum treatment (all superpartners at M\_SUSY)
yields \(\alpha\)\_s = 0.1038 (12.1\% deviation). The physical
\(G_{2}\)-MSSM spectrum is split: squarks and sleptons decouple at
M\_squark = 3165 GeV, while gauginos remain light. Below M\_squark, the
effective theory is SM + gauginos with \(b_{3}\) = \(- 5\) (instead of
\(- 3\) for the full MSSM). This step-function matching gives
\(\alpha\)\_s = 0.1224 (3.7\% deviation), which is the value adopted
throughout.

The RGE deviations reflect sensitivity to the SUSY spectrum,
particularly:

\begin{itemize}

\item
  \textbf{tan \(\beta\) dependence}: The 2-loop terms involving top
  Yukawa are sensitive to tan \(\beta\)
\item
  \textbf{Threshold corrections}: The split-spectrum matching captures
  the dominant effect; residual corrections from the detailed spectrum
  are subdominant
\item
  \textbf{3-loop effects}: Omitted from the current calculation
\end{itemize}

\textbf{Topological vs.~infrared}: The topological \(sin^{2} \theta\)\_W
= 3/13 = 0.23077 (dev 0.19\% from PDG) is more accurate than the
RGE-evolved value (dev 2.78\%). This suggests the topological value may
represent an infrared fixed point rather than a UV boundary condition
: an observation also made in the \(G_{2}\)-MSSM literature \cite{33}.

\textbf{KK threshold corrections}: With 71 KK modes above M\_GUT and
\(\Sigma\) ln(m/M\_GUT) = 89.2, the net correction to
\(\alpha\)\_\(GUT^{-1}\) is 0.0 (smooth \(K_{7}\) manifolds have
vanishing KK threshold corrections due to the cancellation between
towers).

\textbf{SUSY spectrum implications}:

\begin{longtable}[]{@{}
  >{\raggedright\arraybackslash}p{(\linewidth - 4\tabcolsep) * \real{0.4167}}
  >{\raggedright\arraybackslash}p{(\linewidth - 4\tabcolsep) * \real{0.2500}}
  >{\raggedright\arraybackslash}p{(\linewidth - 4\tabcolsep) * \real{0.3333}}@{}}
\toprule\noalign{}
\begin{minipage}[b]{\linewidth}\raggedright
Particle
\end{minipage} & \begin{minipage}[b]{\linewidth}\raggedright
Mass
\end{minipage} & \begin{minipage}[b]{\linewidth}\raggedright
Source
\end{minipage} \\
\midrule\noalign{}
\endhead
\bottomrule\noalign{}
\endlastfoot
m\_gravitino & 166 GeV & F-term from gaugino condensation \\
m\_moduli & 3165 GeV (3.2 TeV) & SU(8) gaugino condensation \\
m\_gluino & 442 GeV & cMSSM with \(m_{1}\)/\textsubscript{2} =
m\_gravitino \\
m\_squark & 1094 GeV & cMSSM (running mass at low scale; decoupling
matched at M\_SUSY = m\_moduli = 3165 GeV) \\
m\_slepton & 175 GeV & cMSSM \\
LSP (Bino) & 70 GeV & Lightest neutralino (pure Bino; viable: LEP
chargino bound 103 GeV does not exclude pure Bino LSP with suppressed
gauge couplings \cite{34}) \\
\end{longtable}


\textbf{Phenomenological caveat}: These masses are computed within the
cMSSM approximation. Standard simplified SUSY searches at ATLAS/CMS
already exclude portions of this parameter space; the spectrum is viable
only in compressed or suppressed-coupling realizations. Definitive
testing requires a dedicated recast of current ATLAS/CMS results against
the \(G_{2}\)-MSSM split-spectrum scenario.

\subsection{The B-Test Identity and Holonomy Sequence}

The gauge coupling predictions \(sin^{2} \theta\)\_W = 3/13 and
\(\alpha\)\_s = $\sqrt{2}$/12, combined with the MSSM beta-function structure,
yield an algebraic identity connecting the fine structure constant to
\(G_{2}\) representation theory.

\textbf{The B parameter}: In any GUT framework, the ``B-test''
quantifies consistency of gauge coupling unification:

\[B = \frac{\alpha_1^{-1} - \alpha_2^{-1}}{\alpha_2^{-1} - \alpha_3^{-1}}\]

For the MSSM with N\_gen = 3 generations and one Higgs doublet pair, the
beta-function coefficients are (\(b_{1}\), \(b_{2}\), \(b_{3}\)) =
(33/5, 1, \(- 3\)), giving B = (\(b_{1} - b_{2}\))/(\(b_{2} - b_{3}\)) =
(28/5)/4 = \textbf{7/5}. The number of generations enters through N\_gen
= \(b_{2}\)/dim(\(K_{7}\)) = 21/7 = 3.

\textbf{Theorem} (B-test): *Given \(sin^{2} \theta\)\_W =
\(b_{2}\)/(\(b_{3}\)+dim(\(G_{2}\))) = 3/13 and \(\alpha\)\_s =
$\sqrt{2}$/(dim(\(G_{2}\))\(- p_{2}\)) = $\sqrt{2}$/12, the MSSM relation B = 7/5 holds
if and only if*

\[\alpha_{\mathrm{em}}^{-1}(M_Z) = (b_3 + \dim(G_2)) \cdot \sqrt{2} = 91\sqrt{2} = 128.693\ldots\]

\emph{where 91 = dim($\Lambda^2 \mathfrak{g}_2$)
is the dimension of the exterior square of the \(G_{2}\) Lie algebra.}

\emph{Proof sketch}. The topological \(sin^{2} \theta\)\emph{W = 3/13
and the GUT normalization factor 3/5 force
\(\alpha\)\textsubscript{1}\textsuperscript{-}\textsuperscript{1} =
2\(\alpha\)\textsubscript{2}\textsuperscript{-}\textsuperscript{1}. Then
B =
\(\alpha\)\textsubscript{2}\textsuperscript{-}\textsuperscript{1}/(\(\alpha\)\textsubscript{2}\textsuperscript{-}\textsuperscript{1}
\(-\)
\(\alpha\)\textsubscript{3}\textsuperscript{-}\textsuperscript{1}), and
B = 7/5 requires \(\alpha\)\textsubscript{2}/\(\alpha\)\textsubscript{3}
= 2/7 = \(p_{2}\)/dim(\(K_{7}\)). Substituting
\(\alpha\)\textsubscript{3}\textsuperscript{-}\textsuperscript{1} =
(dim(\(G_{2}\)) \(-\) \(p_{2}\))/$\sqrt{2}$ = 6$\sqrt{2}$ gives \(\alpha\)}\(em^{-1}\) =
(7 \(\times\) 13)$\sqrt{2}$ = 91$\sqrt{2}$. The factor 91 = \(b_{3}\) + dim(\(G_{2}\)) =
77 + 14 = dim($\Lambda^2 \mathfrak{g}_2$) is a
\(G_{2}\) representation-theoretic invariant. \emph{(Note: the B-test
identity gives the GUT-scale \(\alpha\)\emph{\(em^{-1}\) \(\approx\)
128.7, valid at M\_GUT where B = 7/5; the observed low-energy value
\(\alpha\)}\(em^{-1}\) \(\approx\) 137 is a separate derivation via
§4.1. Both are used consistently in GIFT.)}

\textbf{The holonomy sequence}: At the B = 7/5 exact scale, all three
inverse couplings are integer multiples of $\sqrt{2}$:

\begin{longtable}[]{@{}
  >{\raggedright\arraybackslash}p{(\linewidth - 6\tabcolsep) * \real{0.1887}}
  >{\raggedright\arraybackslash}p{(\linewidth - 6\tabcolsep) * \real{0.1321}}
  >{\raggedright\arraybackslash}p{(\linewidth - 6\tabcolsep) * \real{0.3208}}
  >{\raggedright\arraybackslash}p{(\linewidth - 6\tabcolsep) * \real{0.3585}}@{}}
\toprule\noalign{}
\begin{minipage}[b]{\linewidth}\raggedright
Coupling
\end{minipage} & \begin{minipage}[b]{\linewidth}\raggedright
Value
\end{minipage} & \begin{minipage}[b]{\linewidth}\raggedright
= integer \(\times\) $\sqrt{2}$
\end{minipage} & \begin{minipage}[b]{\linewidth}\raggedright
Topological origin
\end{minipage} \\
\midrule\noalign{}
\endhead
\bottomrule\noalign{}
\endlastfoot
\(\alpha\)\textsubscript{1}\textsuperscript{-}\textsuperscript{1} &
59.40 & 42$\sqrt{2}$ & 2\(b_{2}\) \\
\(\alpha\)\textsubscript{2}\textsuperscript{-}\textsuperscript{1} &
29.70 & 21$\sqrt{2}$ & \(b_{2}\) \\
\(\alpha\)\textsubscript{3}\textsuperscript{-}\textsuperscript{1} & 8.49
& 6$\sqrt{2}$ & dim(\(K_{7}\))\(- 1\) \\
\end{longtable}


Their ratios in lowest terms:

\[\alpha_1^{-1} : \alpha_2^{-1} : \alpha_3^{-1} = \dim(G_2) : \dim(K_7) : p_2 = 14 : 7 : 2\]

This \textbf{holonomy sequence} encodes the \(G_{2}\) structure of
\(K_{7}\) directly in the gauge couplings: the holonomy group dimension,
the manifold dimension, and the Pontryagin ratio \(p_{2}\) =
dim(\(G_{2}\))/dim(\(K_{7}\)).

\textbf{Exact unification}: Using the 1-loop MSSM RGE
\(\alpha\)\emph{\(i^{-1}\)(\(\mu\)) = \(\alpha\)}\(i^{-1}\)(M\_Z) \(-\)
b\_i/(2\(\pi\))\(\cdot\)ln(\(\mu\)/M\_Z), the three couplings unify
exactly at:

\[t_{\mathrm{GUT}} = \frac{15\pi\sqrt{2}}{2}, \qquad \alpha_{\mathrm{GUT}}^{-1} = \frac{\sqrt{2}\,(b_3 - \mathrm{rank}(E_8))}{4} = \frac{69\sqrt{2}}{4} \approx 24.4\]

where 69 = \(b_{3}\) \(-\) rank(\(E_{8}\)) = 77 \(-\) 8 =
\textbar PSL(2,7)\textbar{} \(-\) H* = 168 \(-\) 99. This gives M\_GUT =
M\_Z\(\cdot\)\(e^{t}\) = 2.7 \(\times\) 10\(^{16}\) GeV, consistent with
the standard MSSM estimate. The GUT coupling \(\alpha\)\_\(GUT^{-1}\) =
24.4 compares to the gauge kinetic function value 25.3 used in §5.3
(3.6\% discrepancy, reflecting the approximate nature of the identity).

\textbf{Numerical status}: Using GIFT's topological couplings at M\_Z, B
= 1.4033 (0.23\% from 7/5). This is closer to 7/5 than the purely
experimental value B = 1.3948 (0.37\% off), suggesting the identity has
genuine geometric content. The \textasciitilde0.5\% deficit between 91$\sqrt{2}$
= 128.69 and the GIFT RGE-running prediction
\(\alpha\)\_\(em^{-1}\)(M\_Z) = 128.03 {[}distinct from the topological
low-energy value 137.033 = 4\(\pi\)/\(\alpha\)\_em{]} remains an open
question: it may trace to 2-loop or threshold effects not captured at
the algebraic level.

\subsection{Bundle Universality and Gram Conditioning}

The gauge kinetic function f\_IJ on \(K_{7}\) determines gauge coupling
universality. From the 22 harmonic 2-forms on \(K_{3}\):

\textbf{f\_IJ eigenvalue spectrum}: 22 eigenvalues in {[}0.733,
0.767{]}, with \textbf{cond(f\_IJ) = 1.047}, near-perfect
universality. The gauge coupling ratio \(\alpha\)\_ratio =
\(\alpha\)\_max/\(\alpha\)\_min = 1.000002, confirming that all gauge
couplings are effectively identical at the compactification scale.

\textbf{Gram matrices} quantify orthonormality of the harmonic form
basis:

\begin{longtable}[]{@{}lllll@{}}
\toprule\noalign{}
Matrix & Size & Condition & PD & Off-diag max \\
\midrule\noalign{}
\endhead
\bottomrule\noalign{}
\endlastfoot
G\_K3(22) & 22\(\times 22\) & \textbf{1.05} & Yes & 0.012 \\
G\_K7(22) & 22\(\times 22\) & \textbf{1.05} & Yes & 0.012 \\
G\_35 & 35\(\times 35\) & \textbf{7.66} & Yes &, \\
G\_77 & 77\(\times 77\) & \textbf{7.66} & Yes & 7\(\times 10^{-5}\) \\
\end{longtable}


The K3 and K\textsubscript{7} 22-form bases are nearly orthonormal (condition
\textasciitilde1.05). The full 77-form basis has moderate conditioning
(7.66), with cross-block coupling between constant and fiber modes
bounded by 7 \(\times\) 10\(^{-5}\). Gram-Schmidt orthogonalization
residuals are \textless{} 5 \(\times\) 10\(^{-16}\) (machine precision).

\begin{figure}[htbp]
\centering
\includegraphics[width=0.95\textwidth]{fig_gauge_bundle_v2.png}
\caption{Gauge bundle structure, eigenvalue spectrum of $f_{IJ}$ and Gram matrix diagnostics}
\label{fig4}
\end{figure}

\textbf{Lean certification}: \texttt{TCSGaugeBreaking.lean} (0 axioms,
14 theorems, 10-conjunct master certificate) +
\texttt{GaugeBundleData.lean} (0 axioms, 12 theorems, 11-conjunct master
certificate).

\subsection{Summary}

The gauge sector pipeline covers \(E_{8}\) \(\to\) SM with N\_gen = 3,
all anomalies cancelled, near-perfect bundle universality (cond 1.047),
and Lean-certified results. The topological gauge predictions
(\(sin^{2} \theta\)\_W, \(\alpha\)\_s) are more precise than the
dynamical RGE running, suggesting the topological values may represent
infrared fixed points rather than UV boundary conditions. The B-test
identity (§5.4) reveals that these two predictions, combined with the
MSSM structure, encode the holonomy sequence
dim(\(G_{2}\)):dim(\(K_{7}\)):\(p_{2}\) = 14:7:2 in the gauge coupling
ratios: a direct imprint of \(G_{2}\) geometry on low-energy physics.

\noindent

\rule{\textwidth}{0.2pt}

\section{Mass Hierarchy: From Geometry to Generations}

The five-order-of-magnitude lepton mass hierarchy (m\_e : m\_\(\mu\) :
m\_\(\tau\) \textasciitilde{} 1 : 207 : 3477) is one of the deepest
puzzles in particle physics. This section presents two independent
geometric mechanisms that individually reproduce the hierarchy to
\textasciitilde2--6\% and, when combined, achieve sub-percent precision.

\subsection{\texorpdfstring{Three Generations from the \(Z_{3}\)
Mechanism}{Three Generations from the Z\_\{3\} Mechanism}}\label{three-generations-from-the-z_3-mechanism}

Since \(\pi\)\textsubscript{1}(\(K_{7}\)) = \{1\}, traditional Wilson
line breaking is trivial. Instead, three generations emerge from a
\(Z_{3}\) lattice action on the K3 fiber of the neck region.

\textbf{Wilson line theorem}: The rank of the Wilson line operator is
preserved under perturbation. SVD analysis gives singular values
{[}5.71, 0.62, 2.4 \(\times\) 10\(^{-15}\){]}, confirming rank 3 (the
third eigenvalue is machine zero, consistent with a \(Z_{3}\) quotient
leaving three independent directions).

\textbf{K3 metric properties}: The K3 fiber metric is nearly flat:
conformal range 0.018\% across the fiber, mean anisotropy 1.4\%. This
near-flatness ensures the eigenvalue splitting is controlled by the
fiber geometry rather than by large-scale metric fluctuations.

\subsection{Lepton Mass Hierarchy: Non-Adiabatic Mechanism (wilson\_line)}

The non-adiabatic eigenvalue splitting mechanism operates on the K3
fiber at coupling c = 0.452 and optimized positions {[}0.0, 0.693,
1.400{]}:

\textbf{Eigenvalues}: {[}0.03383, 0.00205, 9.94 \(\times\)
10\(^{-6}\){]}

\begin{longtable}[]{@{}llll@{}}
\toprule\noalign{}
Ratio & GIFT & Experimental & Deviation \\
\midrule\noalign{}
\endhead
\bottomrule\noalign{}
\endlastfoot
m\_\(\tau\)/m\_\(\mu\) & 16.54 & 16.82 & 1.7\% \\
m\_\(\tau\)/m\_e & 3403 & 3477 & 2.1\% \\
m\_\(\mu\)/m\_e & 205.7 & 206.7 & 0.5\% \\
\end{longtable}


The critical coupling c* = \(10^{-3}\)/\textsuperscript{4} = 0.1778
marks the transition between adiabatic (small splitting,
\textasciitilde2 generations) and non-adiabatic (large splitting, 3
generations) regimes. The physical coupling c = 0.452 \textgreater{} c*
places \(K_{7}\) firmly in the three-generation regime.

\textbf{Adiabatic limit}: At c \(\to\) 0, only two generations are
distinguishable (\(m_{1}\)/\(m_{2}\) = 77.6, \(m_{1}\)/\(m_{3}\) \(\to\)
\(\infty\)). The three-generation structure requires c \textgreater{}
c*, which is satisfied by the neck geometry.

\subsection{Instanton Volume Differences (instanton)}

An independent mechanism generates the mass hierarchy from associative
3-cycle volumes. On \(K_{7}\), there are \textbf{57 associative
3-cycles} with volumes in {[}0.00075, 11.109{]}. The mass relation
m\_i/m\_j = exp(\(\Delta\)V) assigns each generation to a cycle.

\textbf{Optimal assignment} (minimizing combined deviation):

\begin{longtable}[]{@{}
  >{\raggedright\arraybackslash}p{(\linewidth - 8\tabcolsep) * \real{0.2727}}
  >{\raggedright\arraybackslash}p{(\linewidth - 8\tabcolsep) * \real{0.1818}}
  >{\raggedright\arraybackslash}p{(\linewidth - 8\tabcolsep) * \real{0.1136}}
  >{\raggedright\arraybackslash}p{(\linewidth - 8\tabcolsep) * \real{0.1818}}
  >{\raggedright\arraybackslash}p{(\linewidth - 8\tabcolsep) * \real{0.2500}}@{}}
\toprule\noalign{}
\begin{minipage}[b]{\linewidth}\raggedright
Assignment
\end{minipage} & \begin{minipage}[b]{\linewidth}\raggedright
Volume
\end{minipage} & \begin{minipage}[b]{\linewidth}\raggedright
\(\Delta\)V
\end{minipage} & \begin{minipage}[b]{\linewidth}\raggedright
Target
\end{minipage} & \begin{minipage}[b]{\linewidth}\raggedright
Deviation
\end{minipage} \\
\midrule\noalign{}
\endhead
\bottomrule\noalign{}
\endlastfoot
V\_e & 11.109 & (&) &, \\
V\_\(\mu\) & 7.838 & \(\Delta\)V(e-\(\mu\)) = 3.271 & ln(16.82) = 2.823
& 15.9\% \\
V\_\(\tau\) & 2.476 & \(\Delta\)V(e-\(\tau\)) = 8.633 & ln(3477) = 8.154
& 5.9\% \\
\end{longtable}


The e-\(\tau\) hierarchy (5 orders of magnitude) is reproduced to 5.9\%,
while the e-\(\mu\) hierarchy (2.3 orders) shows 15.9\% deviation from
the optimal assignment. The total volume range \(\Delta\)V\_range = 8.92
spans the correct order of magnitude.

\begin{figure}[htbp]
\centering
\includegraphics[width=0.95\textwidth]{fig_instanton_hierarchy_v2.png}
\caption{Instanton mass hierarchy, volume spectrum of 57 associative cycles with generation assignments}
\label{fig5}
\end{figure}

\subsection{Combined wilson\_line+instanton Pipeline}

The key insight is that wilson\_line and instanton probe different
aspects of \(K_{7}\) geometry: - \textbf{wilson\_line}: Fiber geometry
\(\to\) eigenvalue spacing (relative structure) - \textbf{instanton}:
Cycle volumes \(\to\) exponential hierarchy (absolute scale)

The two mechanisms are connected by \(\alpha\) = \(e^K\) =
exp(\(K_{0}\)) = $\hat{V}^{-3}$, where \(K_{0}\) = \(- 5\).891 is the
Kähler potential of \(K_{7}\) (§6.5). This is a purely geometric
quantity: the instanton action normalization derived from the
compactification volume: not a fit parameter:

\begin{longtable}[]{@{}
  >{\raggedright\arraybackslash}p{(\linewidth - 10\tabcolsep) * \real{0.1296}}
  >{\raggedright\arraybackslash}p{(\linewidth - 10\tabcolsep) * \real{0.1667}}
  >{\raggedright\arraybackslash}p{(\linewidth - 10\tabcolsep) * \real{0.1667}}
  >{\raggedright\arraybackslash}p{(\linewidth - 10\tabcolsep) * \real{0.1852}}
  >{\raggedright\arraybackslash}p{(\linewidth - 10\tabcolsep) * \real{0.2407}}
  >{\raggedright\arraybackslash}p{(\linewidth - 10\tabcolsep) * \real{0.1111}}@{}}
\toprule\noalign{}
\begin{minipage}[b]{\linewidth}\raggedright
Ratio
\end{minipage} & \begin{minipage}[b]{\linewidth}\raggedright
wilson\_line raw
\end{minipage} & \begin{minipage}[b]{\linewidth}\raggedright
instanton raw
\end{minipage} & \begin{minipage}[b]{\linewidth}\raggedright
Combined (\(\alpha\) = \(e^K\))
\end{minipage} & \begin{minipage}[b]{\linewidth}\raggedright
Experimental
\end{minipage} & \begin{minipage}[b]{\linewidth}\raggedright
Dev
\end{minipage} \\
\midrule\noalign{}
\endhead
\bottomrule\noalign{}
\endlastfoot
m\_\(\tau\)/m\_e & 3403 (2.1\%) & exp(8.63) & \textbf{3485} & 3477 &
\textbf{0.24\%} \\
m\_\(\tau\)/m\_\(\mu\) & 16.54 (1.7\%) & exp(3.27) & \textbf{16.69} &
16.82 & \textbf{0.75\%} \\
m\_\(\mu\)/m\_e & 205.7 (0.5\%) &, & \textbf{208.8} & 206.8 &
\textbf{0.97\%} \\
\end{longtable}


All three ratios within 1\%: a significant improvement over the
individual mechanisms. The combined pipeline works because the
mechanisms are \textbf{complementary}: wilson\_line provides fine
structure from the eigenvalue splitting, instanton provides the overall
exponential scale from cycle volumes. The key insight is that \(\alpha\)
= \(e^K\) has a natural M-theory interpretation: for M2-branes wrapping
associative 3-cycles, the instanton action scales as S\_inst = \(e^K\)
\(\times\) Vol(\(\Sigma\)) = Vol(\(K_{7}\))\^{}\{\(- 3\)\} \(\times\)
Vol(\(\Sigma\)), giving the correct suppression without any free
parameters.

\textbf{Lean certification}: \texttt{AssociativeVolumes.lean} (0 axioms,
19 theorems, 14-conjunct master certificate) certifies the combined
wilson\_line+instanton results including all three mass ratios and the
geometric \(\alpha\) = \(e^K\).

\subsection{4D Effective Theory (S9)}

Dimensional reduction of the \(G_{2}\) metric yields an N = 1
supergravity theory in 4D:

\textbf{Particle content}: From Betti numbers directly: - 1 gravity
multiplet - 21 vector multiplets (= \(b_{2}\)) - 77 chiral multiplets (=
\(b_{3}\)) - Total: 99 = H* (the effective cohomological dimension)

\textbf{Kähler potential}: \(K_{0}\) = \(- 5\).891, with \(e^K\) =
0.002763 = $\hat{V}^{-3}$. This geometric quantity serves as the
instanton action normalization \(\alpha\) = \(e^K\) in the combined
wilson\_line+instanton pipeline (§6.4). The 7D volume $\hat{V} = 7.126$. The
metric determinant det(\(g_{7}\)) = 65/32 = 2.03125 is locked to its
topological value.

\textbf{Moduli metric}: Total dimension 79 (= 77 physical + 2 gauge),
rank 77 (2 null directions from gauge redundancy). The condition number
cond = 7.66 measures the ratio of largest to smallest moduli mass. The
35 constant-block eigenvalues span {[}1.66, 12.69{]}; the 42 fiber-block
eigenvalues cluster around 6.1.

\textbf{Gaugino condensation}: The SU(8) hidden sector is preferred:

\begin{longtable}[]{@{}llll@{}}
\toprule\noalign{}
Condensate & \(b_{0}\) & \(\Lambda\) (GeV) & m\_moduli (GeV) \\
\midrule\noalign{}
\endhead
\bottomrule\noalign{}
\endlastfoot
SU(3) & 9 & 4.3 \(\times\) 10\(^{8}\) & 1.3 \(\times\) 10\(^{-11}\) \\
SU(5) & 15 & 5.0 \(\times\) 10\(^{11}\) & 0.021 \\
\textbf{SU(8)} & \textbf{24} & \textbf{2.66 \(\times\) 10\(^{13}\)} &
\textbf{3165} \\
\end{longtable}


The SU(8) sector gives m\_moduli = 3165 GeV (3.2 TeV), m\_gravitino =
166 GeV, consistent with the \(G_{2}\)-MSSM spectrum.

\textbf{Physical predictions from the effective theory}: - Proton
lifetime: \(\tau\)\_p = 4.06 \(\times\) 10\(^{38}\) years (well above
current bound \textasciitilde{} 10\(^{34}\) years) - Yukawa effective
rank = 3 (three massive generations) - \(\lambda\)\_physical = 1.358
(physical Yukawa coupling)

\begin{figure}[htbp]
\centering
\includegraphics[width=0.95\textwidth]{fig_kk_spectrum_v2.png}
\caption{KK spectrum: Kaluza-Klein tower structure and 7D Weyl law}
\label{fig6}
\end{figure}

\noindent

\rule{\textwidth}{0.2pt}

\section{Sensitivity Analysis}

A framework claiming 95 observables from 20 structural constants must
address the question: how constrained are these predictions? This
section presents five independent analyses demonstrating that the
predictions are overdetermined, cross-coupled, and statistically
incompatible with coincidence.

\subsection{Formula Structure Analysis}

\textbf{Question}: Do more complex formulas systematically produce
smaller deviations? If so, the framework might be ``fitting'' by formula
complexity.

\textbf{Method}: Random Forest regression with leave-one-out
cross-validation, predicting \textbar deviation\textbar{} from formula
features (number of constants used, maximum constant value, arithmetic
operations, sector membership).

\textbf{Result}: \(R^{2}\) = \(- 0\).518 (LOO-CV), worse than a mean
predictor. The top feature importance is max\_constant\_value (0.41),
followed by n\_expressions (0.13) and sector\_PMNS (0.12). Formula
complexity does \textbf{not} predict deviation magnitude.

\textbf{Verdict}: This is inconsistent with systematic cherry-picking.
Complex formulas do not perform better than simple ones; the precision
is distributed uniformly across formula types.

\subsection{Topological Constant Sensitivity}

\textbf{Method}: Perturb each of the 20 structural constants by
\(\pm 1\) \(\to\) 20 \(\times\) 33 sensitivity matrix S\_ij =
\(\partial\)(observable\_i)/\(\partial\)(constant\_j).

\textbf{SVD analysis}: 19 significant singular values (1 zero,
corresponding to a redundant constant combination). The singular value
spectrum decays smoothly: {[}3.50, 2.54, 2.39, 2.14, 2.11, 2.01, 1.93,
1.76, 1.68, 1.63, \ldots{]}.

\textbf{Sensitivity-deviation correlation}: \(\rho\) = \(- 0\).083,
\textbf{no systematic pattern}. Observables with high sensitivity to
constant perturbations do not have larger deviations. The most sensitive
constants are dim(\(K_{7}\)) (0.068), dim(\(G_{2}\)) (0.062), and H*
(0.062).

\textbf{NK ball rigidity}: The Newton-Kantorovich certification
establishes \(\delta\)g/g \(\leq\) 1.35 \(\times\) 10\(^{-7}\),
confirming the metric is effectively unique within its basin. The
coefficient of variation of metric eigenvalues is \textasciitilde{}
10\(^{-7}\), confirming rigidity.

\subsection{Effective Degrees of Freedom}

\textbf{Method}: SVD of the 20 \(\times\) 33 constant-usage Jacobian
(binary: 1 if constant appears in formula, 0 otherwise). The effective
rank is defined by the singular value decay profile.

\textbf{Results}: - \textbf{r\_eff = 15.53} effective parameters (out of
20 structural constants) - 12 singular values capture 90\% of variance -
17 singular values capture 99\% of variance - 1 zero singular value
(exact linear dependence)

\textbf{Overdetermination ratio}: 33 observables / 15.53 effective
parameters = \textbf{2.13\(\times\)}. The system has more than twice as
many constraints as degrees of freedom: a hallmark of an
overdetermined (not fitted) system.

\textbf{Most-used constants}: \(b_{2}\) appears in 9 observables,
dim(\(G_{2}\)) in 7, dim(\(E_{8}\)) in 7, Weyl in 6, H* in 6. No
constant is used in more than 27\% of observables, confirming broad
coverage rather than concentration.

\begin{figure}[htbp]
\centering
\includegraphics[width=0.95\textwidth]{fig_constant_usage_v2.png}
\caption{Constant usage: SVD spectrum showing effective rank and per-constant usage frequency}
\label{fig7}
\end{figure}

\subsection{Cross-Correlations}

\textbf{Jaccard similarity}: Of the C(33,2) = 528 observable pairs,
\textbf{155 share at least one structural constant} (29.4\% coupled).
All 33 observables belong to a single connected component: no
prediction is isolated.

\textbf{Strong correlations}: 51 pairs have \textbar{}\(\rho\)\textbar{}
\(\geq\) 0.5. Notable examples: - m\_b/m\_t vs m\_\(\mu\)/m\_\(\tau\):
\(\rho\) = \(- 0\).816 (both involve 2\(b_{2}\) = 42) -
\(\alpha\)\textsuperscript{-}\textsuperscript{1} vs
\(\Omega\)\_DM/\(\Omega\)\_b: \(\rho\) = \(- 0\).721 (both involve
rank(\(E_{8}\))) - \(sin^{2} \theta\)\_W vs Q\_Koide: \(\rho\) = +0.673
(both involve \(b_{2}\) and dim(\(G_{2}\)))

\textbf{Mean sector Jaccard} = 0.293: sectors share \textasciitilde29\%
of structural constants, creating a web of inter-sector constraints.

\begin{figure}[htbp]
\centering
\includegraphics[width=0.95\textwidth]{fig_sensitivity_heatmap_v2.png}
\caption{Sensitivity heatmap, constant-observable coupling matrix showing cross-sector correlations}
\label{fig8}
\end{figure}

\begin{figure}[htbp]
\centering
\includegraphics[width=0.95\textwidth]{fig_observable_correlations_v2.png}
\caption{Observable correlations, pairwise Pearson correlations between observable deviations}
\label{fig9}
\end{figure}

\subsection{Coincidence Test}

\textbf{Question}: What is the probability that 33 independent random
predictions match experiment as well as GIFT?

\subsubsection{Uniform Null Model}

Multiple statistical tests under the null hypothesis of uniform
deviations in {[}0, 50\%{]}:

\begin{longtable}[]{@{}llll@{}}
\toprule\noalign{}
Test & Statistic & df & p-value \\
\midrule\noalign{}
\endhead
\bottomrule\noalign{}
\endlastfoot
\(\chi\)\textsuperscript{2} & 1063 & 33 & 5.0 \(\times\)
10\(^{-202}\) \\
Fisher combined & 1100.5 & 66 & 2.2 \(\times\) 10\(^{-187}\) \\
KS (pull normality) & 0.189 &, & 0.165 \\
\textbf{Combined uniform} & (&) & \textbf{\(10^{-346}\)} \\
\end{longtable}


\textbf{Pull distribution}: mean = \(- 0\).774, std = 5.62. 72.7\% of
pulls within 1\(\sigma\), 87.9\% within 2\(\sigma\). The KS test (p =
0.165) is consistent with Gaussian pulls: the deviations are not
cherry-picked.

\textbf{Reduced \(\chi\)\textsuperscript{2}} = 32.2 is large, driven by
outliers (\(\delta\)\_CP, \(\alpha\)\_s(RGE)). The Bayes factor
lo\(g_{10}\) = \(- 2\).3 is inconclusive, reflecting this: the bulk of
predictions match extremely well, but a few outliers inflate the tails.
Removing the 4 known outliers gives reduced \(\chi\)\textsuperscript{2}
\textless{} 2.

\subsubsection{Algebraic Null Model (algebraic\_MC)}

The uniform null model assumes predictions are random numbers. A
stronger test: what if we use random \emph{algebraic formulas} from the
same 20 structural constants? We generated 4,188,086 unique formula
values via exhaustive depth-1/2 enumeration (1.8M) plus 3M random
expression trees (depth 2--4), using 5 binary operations (+, \(-\),
\(\times\), \(\div\), \^{}) and 5 unary transforms (id, inv, $\sqrt{\cdot}$,
\textsuperscript{2}, ln).

\begin{longtable}[]{@{}
  >{\raggedright\arraybackslash}p{(\linewidth - 6\tabcolsep) * \real{0.1633}}
  >{\raggedright\arraybackslash}p{(\linewidth - 6\tabcolsep) * \real{0.1224}}
  >{\raggedright\arraybackslash}p{(\linewidth - 6\tabcolsep) * \real{0.5510}}
  >{\raggedright\arraybackslash}p{(\linewidth - 6\tabcolsep) * \real{0.1633}}@{}}
\toprule\noalign{}
\begin{minipage}[b]{\linewidth}\raggedright
Metric
\end{minipage} & \begin{minipage}[b]{\linewidth}\raggedright
GIFT
\end{minipage} & \begin{minipage}[b]{\linewidth}\raggedright
Random algebraic (median)
\end{minipage} & \begin{minipage}[b]{\linewidth}\raggedright
Factor
\end{minipage} \\
\midrule\noalign{}
\endhead
\bottomrule\noalign{}
\endlastfoot
Mean deviation & \textbf{0.73\%} & 4.1 \(\times\) 10\(^{9}\) \% & 5.6
\(\times\) 10\(^{9}\) \\
Exact matches (\textless0.01\%) & \textbf{5} & 0 &, \\
Within 1\% & \textbf{28} & 0 &, \\
\end{longtable}


\textbf{Per-observable}: combined P(random algebraic formula matches
GIFT's precision) = \textbf{\(10^{-133}\)} (product over 33
observables). Only 0.02\% of 4.2M formulas match \emph{any} observable
within 0.01\%.

\textbf{Set-level}: 0 out of 3,000,000 random sets of 33 algebraic
formulas achieved GIFT's mean deviation, exact count, or within-1\%
count. Under this null: P \textless{} 3.3 \(\times\) 10\(^{-7}\) for
each metric.

\begin{figure}[htbp]
\centering
\includegraphics[width=0.95\textwidth]{fig_mc_per_observable.png}
\caption{Per-observable lo$g_{10}$ P(random match)}
\label{fig10}
\end{figure}

\textbf{Interpretation}: The algebraic null model is far more generous
than the uniform null (it generates formulas from the \emph{same
constants}), yet GIFT's performance remains unmatched across 3M+ trials.
Under both declared null models, chance agreement is excluded at extreme
significance: P = \(10^{-346}\) (uniform) and P = \(10^{-133}\)
(algebraic). These figures quantify agreement under each declared model;
they do not constitute an absolute probability that the framework is correct.

\subsection{PSLQ Residual Analysis (PSLQ\_residual)}

Beyond the statistical null models, we apply PSLQ integer relation
detection to the relative residuals r\_i = (GIFT\_i - exp\_i)/exp\_i of
the 33 Type I observables with experimental comparison. The goal is to
identify whether deviations have structural content, i.e., whether
residuals match algebraic expressions built from the same topological
constants.

\textbf{Method}: For each observable, we compute the residual r\_i and
test it against: (1) rational approximations p/q with small
denominators, (2) structural fractions involving GIFT constants
(dim(\(G_{2}\)), dim(\(E_{8}\)), \(b_{2}\), \(b_{3}\), H*, PSL(2,7),
etc.), (3) PSLQ integer relations with the 20 structural constants, (4)
mpmath.identify() for closed-form recognition.

\textbf{Key findings}:

\begin{longtable}[]{@{}
  >{\raggedright\arraybackslash}p{(\linewidth - 6\tabcolsep) * \real{0.2500}}
  >{\raggedright\arraybackslash}p{(\linewidth - 6\tabcolsep) * \real{0.2292}}
  >{\raggedright\arraybackslash}p{(\linewidth - 6\tabcolsep) * \real{0.2500}}
  >{\raggedright\arraybackslash}p{(\linewidth - 6\tabcolsep) * \real{0.2708}}@{}}
\toprule\noalign{}
\begin{minipage}[b]{\linewidth}\raggedright
Observable
\end{minipage} & \begin{minipage}[b]{\linewidth}\raggedright
Residual r
\end{minipage} & \begin{minipage}[b]{\linewidth}\raggedright
Best match
\end{minipage} & \begin{minipage}[b]{\linewidth}\raggedright
Match error
\end{minipage} \\
\midrule\noalign{}
\endhead
\bottomrule\noalign{}
\endlastfoot
\(\delta\)\_CP & +0.1130 & dim(\(G_{2}\))/(dim(\(E_{8}\))/2) = 14/124 &
0.08\% \\
m\_b/m\_t & -0.0079 & -1/(N\_gen \(\times\) 2\(b_{2}\)) = -1/126 &
exact \\
Y\_p & +0.00368 & 1/(\(\varphi\) \(\times\) PS\(L_{27}\)) =
1/(\(\varphi\) \(\times\) 168) & 0.04\% \\
\(sin^{2} \theta\)\textsubscript{1}\textsubscript{2}(CKM) & +0.00372 &
6/(\(b_{2}\) \(\times\) \(b_{3}\)) = 6/1617 & 0.2\% \\
\end{longtable}


\textbf{Documented structural observation}: Only \(\delta\)\_CP. A possible compactification factor (62/69 = dim(\(E_{8}\))/(dim(\(E_{8}\)) + 4 dim(\(K_{7}\)))) is recorded as a post-hoc observation in S2 Appendix F; it is not adopted as a revision and the canonical GIFT prediction remains 197\({}^{\circ}\).

\textbf{Not adopted}: The m\_b/m\_t correction (-1/126) and Y\_p
correction (1/(\(\varphi\) \(\times\) 168)) are documented for future
work pending structural derivation from the compactification geometry.
The CKM correction (6/(\(b_{2}\) \(\times\) \(b_{3}\))) is suggestive
but at lower precision.

\textbf{Epistemological note}: GIFT evolved iteratively from 6 free
parameters (v1) through 4, 3, and finally 0 (v3). Each refinement added
structural content while removing degrees of freedom. The \(\delta\)\_CP
compactification factor is documented because it has geometric meaning
independent of the experimental data it happens to match, but is not
adopted as a revision: the raw topological 197\({}^{\circ}\) stands
as the prediction.

\noindent

\rule{\textwidth}{0.2pt}

\section{Formal Verification and Statistical Analysis}

\subsection{Lean 4 Verification}

The GIFT framework is formally verified in Lean 4 \cite{28} with Mathlib
\cite{29}:

\begin{longtable}[]{@{}ll@{}}
\toprule\noalign{}
Category & Count \\
\midrule\noalign{}
\endhead
\bottomrule\noalign{}
\endlastfoot
Source files & 134 (128 core + 6 generated; 12 test/support) \\
Build jobs & 8378 \\
Unproven (sorry) & 0 \\
Published axioms & 4 (all substantive) \\
Certificate conjuncts & \textbf{213} \\
\end{longtable}


The conjuncts cover metric/torsion/topology (38),
couplings/masses/mixing (55), KK spectrum (41), metric eigenvalues (15),
spectral invariants (10), \(\delta\)\_CP compactification (6), gauge
breaking (10), bundle universality (11), instanton hierarchy (14), and
7D Weyl law (7). Full per-file breakdown in Supplement S3, §3.7.

\begin{verbatim}
theorem weinberg_relation :
  b2 * 13 = 3 * (b3 + dim_G2) := by native_decide

theorem koide_relation :
  dim_G2 * 3 = b2 * 2 := by native_decide
\end{verbatim}

The \(E_{8}\) root system is fully proven (12/12 theorems, basis
generation). \(G_{2}\) differential geometry (exterior algebra, Hodge
star, torsion-free condition) is axiom-free. \(G_{2}\) group structure
(\texttt{g2\_mul\_closed}, \texttt{g2\_subset\_SO7},
\texttt{g2\_det\_mul\_gram}) is proven by \texttt{native\_decide}. The 4 substantive axioms cover: \texttt{K7\_analysis\_data}
(HarmonicForms.lean), \texttt{K7\_spectral\_data} (SpectralTheory.lean),
\texttt{literature\_package} (LiteratureAxioms.lean), and
\texttt{KK\_YM\_EFT} (KKSpectralBridge.lean).

Three certified results anchor the formalization: (1) \textbf{\(G_{2}\)
three-form}: Bryant-Joyce \(\varphi\)\textsubscript{0} on
\(\mathbb{R}\)\textsuperscript{7} formalized in
\texttt{G2ThreeForm.lean}, all 7 nonzero coefficients certified by
\texttt{native\_decide}, dim(\(g_{2}\)) = 14 proven; (2)
\textbf{$\bar{\nu}(K_{7}, g) = 0$}: CGN invariant certified zero via
rectangular TCS ($k_+ = k_- = 1$, CGN Main Corollary \cite{16}); (3)
\textbf{KK spectral bridge}, 4D mass gap formally conditional on
KK\_YM\_EFT alone, all spectral ingredients Lean-certified.

\subsection{Observable Coverage}

Of the 95 observables, \textbf{55 are Lean-certified}:

\begin{longtable}[]{@{}llll@{}}
\toprule\noalign{}
Type & Certified & Total & Coverage \\
\midrule\noalign{}
\endhead
\bottomrule\noalign{}
\endlastfoot
\textbf{I} & 33 & 33 & 100\% \\
\textbf{II} & 0 & 19 & 0\% \\
\textbf{III} & 14 & 21 & 67\% \\
\textbf{IV} & 8 & 22 & 36\% \\
\textbf{Total} & \textbf{55} & \textbf{95} & \textbf{58\%} \\
\end{longtable}


Type II observables are Type I ratios \(\times\) experimental VEVs
(e.g., m\_u = m\_u/m\_d \(\times\) m\_d(PDG)). The algebraic step,
the ratio, is Lean-certified for all 33 core Type I formulas; only
the physical scale identification step is uncertified. Axiomatizing VEV
inputs would be circular (they are experimental inputs, not
predictions). Type III coverage includes the new gauge (10+11 conjuncts)
and instanton (14 conjuncts) certificates.

\subsection{Scope of Verification}

\textbf{What is proven}: Arithmetic identities relating topological
integers. Given \(b_{2}\) = 21, \(b_{3}\) = 77, dim(\(G_{2}\)) = 14,
etc., the numerical relations are machine-verified.

\textbf{What is not proven}: Existence of \(K_{7}\), physical
interpretation of ratios as SM parameters, uniqueness of formula
assignments. The verification establishes \textbf{internal consistency},
not physical truth.

\subsection{Statistical Uniqueness}

Among 192,349 alternative (gauge group, holonomy, Betti) configurations
tested by Monte Carlo, zero outperform GIFT: mean deviation 0.73\% vs
32.9\% for alternatives (P \textless{} 5 \(\times\) 10\(^{-6}\),
\textgreater{} 4.5\(\sigma\)). \(E_{8} \times E_{8}\) beats all tested
groups by 13\(\times\); \(G_{2}\) holonomy beats SU(3) (Calabi-Yau) by
13\(\times\). Only rank 8 gives N\_gen = 3 exactly. The sensitivity
analysis of §7 provides complementary evidence: (1) parameter variation
(P \textless{} 5 \(\times\) 10\(^{-6}\)), (2) uniform coincidence (P =
10\(^{-346}\)), and (3) algebraic coincidence (4.2M random formulas, P =
10\(^{-133}\)). Full gauge-group and holonomy rankings in Supplement S2,
§23.

\subsection{\texorpdfstring{The \(G_{2}\)
Metric}{The G\_\{2\} Metric}}\label{the-g_2-metric}

The predictions in §4 depend only on topological invariants. However,
the \(G_{2}\) metric constrained by det(g) = 65/32 is numerically
constructed as a Chebyshev-Cholesky expansion with 169 parameters
(companion paper [A]):

\begin{longtable}[]{@{}
  >{\raggedright\arraybackslash}p{(\linewidth - 2\tabcolsep) * \real{0.5882}}
  >{\raggedright\arraybackslash}p{(\linewidth - 2\tabcolsep) * \real{0.4118}}@{}}
\toprule\noalign{}
\begin{minipage}[b]{\linewidth}\raggedright
Quantity
\end{minipage} & \begin{minipage}[b]{\linewidth}\raggedright
Value
\end{minipage} \\
\midrule\noalign{}
\endhead
\bottomrule\noalign{}
\endlastfoot
‖T‖\_\(C^{0}\) (torsion) & 2.949 \(\times\) 10\(^{-5}\) \\
NK parameter h (analytical, [A]) & 8.95 \(\times\)
10\(^{-9}\) (\(\beta\) = 0.321 exact, margin \(\times 56\)M) \\
NK parameter h (numerical, tighter) & 1.43 \(\times\) 10\(^{-9}\)
(\(\beta\) = 0.0296 numerical, margin \(\times 350\)M) \\
\(\delta\)g/g (analytical ball radius) & \(\leq\) 4.86 \(\times\)
10\(^{-6}\) \\
det(g) & 65/32 (exact, by construction) \\
\end{longtable}


The interval-arithmetic Newton-Kantorovich certification [A]
establishes existence and uniqueness of a torsion-free \(G_{2}\) metric
within the certified ball. The analytical certificate uses \(\beta\) =
0.321 derived from det(g) = 65/32 (exact constraint); a tighter
numerical estimate \(\beta\) = 0.0296 yields h = 1.43 \(\times\)
10\(^{-9}\) (margin \(\times 350\)M) but rests on numerical Lipschitz
estimates rather than the analytical bound. Both values certify h
\textless{} 0.5 (Kantorovich threshold) by a factor of at least 56
million. Torsion reduction \(\times 2995\) in 5 Joyce iterations is well
within Joyce's perturbative regime.

\textbf{Analytic invariant.} The Crowley-Goette-Nordström invariant
$\bar{\nu}(K_{7}, g) \in \mathbb{Z}$ \cite{16} vanishes for any
rectangular TCS with twisting numbers $k_+ = k_- = 1$, by CGN Main Corollary
(gluing angle \(\theta\) = \(\pi\)/2 forced). This conditional result is
certified in Lean (\texttt{TCSConstruction.lean},
\texttt{K7\_nu\_bar\_zero}): if \(K_{7}\) is realized as a rectangular
TCS, then $\bar{\nu}(K_{7}, g) = 0$ without any additional geometric
computation. The building block identification remains open.

\textbf{Results from the NK-certified metric.} Analysis of the
NK-certified metric yields several structural results:

\begin{itemize}

\item
  \textbf{V\_min formula}: The minimum associative cycle volume
  satisfies $V_{\min} = \sqrt{\mathrm{Vol}(K_7)/11}$, where 11 = \(b_{3}\)/n =
  77/7. NK numerical value 219.90; formula gives 221.24 (0.6\%
  agreement).
\item
  \textbf{Harmonic decompositions}: \(b_{2}\) = 7+14 = 3+18 = 11+10
  (\(G_{2}\) reps / hyperkähler triple / TCS blocks); \(b_{3}\) = 35+42
  = (1+7+27)+2\(\times 21\); spectral gap 10522\(\times\) between zero
  and non-zero modes.
\item
  \textbf{U(1)\textsuperscript{2} exact symmetry}: Period integrals
  S\_\(\theta\) = S\_\(\psi\) = 6.1265, exact to 2.6\(\times 10^{-8}\)
, propagates from metric to all period integrals.
\item
  \textbf{Universal law}: \(\lambda\)\textsubscript{1} \(\times\) H* =
  12.3364 holds for all 66 known \(G_{2}\) manifolds.
\item
  \textbf{Lepton hierarchy from periods}: ln(m\_\(\tau\)/m\_e) = 8.154
  from SD associative volumes \(\to\) \(e^{8.154}\) = 3477
  (Lean-certified).
\end{itemize}

Full details in the companion paper [A].

\noindent

\rule{\textwidth}{0.2pt}

\section{Discussion, Falsifiability, and Conclusion}

\subsection{Falsifiable Predictions}

The framework makes concrete, testable predictions. The most critical:

**\(\delta\)\_CP**: The topological prediction is \(\delta\)\_CP =
197\({}^{\circ}\) = 7 \(\times\) 14 + 99 (pure geometry). NuFIT 6.0 (NO
w/o SK) gives 177\({}^{\circ}\) \(\pm\) 20\({}^{\circ}\), an 11.3\%
deviation from center but at the edge of the 1\(\sigma\) uncertainty
band. The experimental central value has shifted significantly between
NuFIT releases (197\({}^{\circ}\) in 5.2, 177\({}^{\circ}\) in 6.0) and
may shift further as DUNE accumulates statistics (2028--2040). A possible compactification factor (62/69) is recorded as a post-hoc observation in S2 Appendix F and is not part of the main prediction. \textbf{Falsification}: \(\delta\)\_CP outside
{[}182, 212{]}\({}^{\circ}\) at 3\(\sigma\) creates serious tension.

\textbf{N\_gen = 3}: No flexibility. A fourth generation immediately
falsifies.

\textbf{m\_s/m\_d = 20}: Lattice QCD target precision \(\pm 0\).5 by
2030. Current 20.0 \(\pm\) 1.0.

**\(sin^{2} \theta\)\_W = 3/13**: FCC-ee precision \textasciitilde{}
10\(^{-5}\) (4\(\times\) improvement over current).

\textbf{New tests from Types II/III}: - Combined lepton ratios
(wilson\_line+instanton): \(\tau\)/e = 3485, \(\tau\)/\(\mu\) = 16.69,
\(\mu\)/e = 208.8, all within 1\% - Proton lifetime: \(\tau\)\_p =
4.06 \(\times\) 10\(^{38}\) years (beyond near-term Hyper-K sensitivity;
Hyper-K constrains nearby GUT alternatives) - SUSY spectrum:
m\_gravitino = 166 GeV, m\_moduli = 3.2 TeV (model-dependent; viable
only in compressed/suppressed-coupling realizations, requiring dedicated
recast at HL-LHC)

\begin{longtable}[]{@{}llll@{}}
\toprule\noalign{}
Experiment & Observable & Timeline & Test Level \\
\midrule\noalign{}
\endhead
\bottomrule\noalign{}
\endlastfoot
DUNE Phase I & \(\delta\)\_CP (3\(\sigma\)) & 2028--2030 & Critical \\
DUNE Phase II & \(\delta\)\_CP (5\(\sigma\)) & 2030--2040 &
Definitive \\
Lattice QCD & m\_s/m\_d & 2028--2030 & Strong \\
HL-LHC & SUSY spectrum & 2029--2040 & Complementary \\
Hyper-Kamiokande & \(\delta\)\_CP, \(\tau\)\_p & 2034+ &
Complementary \\
FCC-ee & \(sin^{2} \theta\)\_W, Q\_Koide & 2040s & Definitive \\
\end{longtable}


\subsection{Relation to M-Theory}

\(E_{8} \times E_{8}\) and \(G_{2}\) holonomy connect directly to
M-theory {[}33,34{]}: - Heterotic string theory requires
\(E_{8} \times E_{8}\) for anomaly cancellation \cite{19} - M-theory on
\(G_{2}\) manifolds preserves N = 1 SUSY in 4D \cite{34}

GIFT differs from standard M-theory phenomenology \cite{35} by focusing
on topological invariants rather than moduli stabilization. Where
M-theory faces the landscape problem (\textasciitilde{}\(10^{500}\)
vacua), GIFT proposes that topological data alone constrain the physics.
The \(G_{2}\)-MSSM spectrum (§5.3, §6.5) is consistent with the
phenomenology of Acharya et al.~\cite{33}: m\_gravitino \textasciitilde{}
100 GeV, m\_moduli \textasciitilde{} few TeV, gaugino condensation in a
hidden sector.

The second \(E_{8}\) factor is required by anomaly cancellation but has
no direct coupling to Standard Model fields. In heterotic M-theory,
gaugino condensation in this hidden sector drives SUSY breaking
\cite{33}. A natural interpretation is that the hidden \(E_{8}\) sector
provides the dark sector: the predicted cosmological ratios
\(\Omega\)\_DM/\(\Omega\)\_b = 43/8 and \(\Omega\)\_DE = ln(2)
\(\times\) 98/99 emerge from the same topological invariants (\(b_{2}\),
\(b_{3}\)) that also determine dim(\(E_{8}\)) = 248. Whether this
connection is structural or coincidental remains an open question.

\subsection{Comparison}

\begin{longtable}[]{@{}
  >{\raggedright\arraybackslash}p{(\linewidth - 6\tabcolsep) * \real{0.2500}}
  >{\raggedright\arraybackslash}p{(\linewidth - 6\tabcolsep) * \real{0.1364}}
  >{\raggedright\arraybackslash}p{(\linewidth - 6\tabcolsep) * \real{0.4091}}
  >{\raggedright\arraybackslash}p{(\linewidth - 6\tabcolsep) * \real{0.2045}}@{}}
\toprule\noalign{}
\begin{minipage}[b]{\linewidth}\raggedright
Criterion
\end{minipage} & \begin{minipage}[b]{\linewidth}\raggedright
GIFT
\end{minipage} & \begin{minipage}[b]{\linewidth}\raggedright
String Landscape
\end{minipage} & \begin{minipage}[b]{\linewidth}\raggedright
Lisi \(E_{8}\)
\end{minipage} \\
\midrule\noalign{}
\endhead
\bottomrule\noalign{}
\endlastfoot
Falsifiable predictions & Yes (\(\delta\)\_CP = 197\({}^{\circ}\),
N\_gen) & Not yet (landscape selection) & Not yet (embedding
obstruction) \\
Adjustable parameters & 0 & \textasciitilde{}\(10^{500}\) & 0 \\
Formal verification & 213 conjuncts, 4 axioms & No & No \\
Precise predictions & 95 (66 testable) & Qualitative &
\textasciitilde5 \\
Gauge breaking & \(E_{8} \to\)SM (6 levels) & Landscape-dependent &
Single \(E_{8}\) \\
Mass hierarchy & 2 mechanisms + combined & (&) \\
Sensitivity analysis & r\_eff = 15.53, P = \(10^{-346}\) & (&) \\
\end{longtable}


\textbf{Distler-Garibaldi obstruction} \cite{36}: Lisi's \(E_{8}\) ToE
attempted to embed all Standard Model fermions (including chirality)
directly as roots of a single \(E_{8}\). Distler and Garibaldi proved
this is mathematically impossible: no \(E_{8}\) representation
decomposes into the correct chiral SM spectrum. GIFT avoids this
obstruction entirely: \(E_{8} \times E_{8}\) is the gauge group of
heterotic M-theory (not a particle container), the SM gauge group
emerges from a six-level breaking chain (§5.1), fermion generations are
fixed by the topological constraint (rank(\(E_{8}\)) + N\_gen)\(b_{2}\)
= N\_gen\(\cdot b_{3}\), giving N\_gen = 3 (§2), and chirality is a
topological property of the compactification. The mathematical
relationship between \(E_{8}\) and particles is cohomological (emergence
from geometry), not representational (embedding into algebra).

\subsection{Limitations and Open Questions}

\begin{longtable}[]{@{}
  >{\raggedright\arraybackslash}p{(\linewidth - 4\tabcolsep) * \real{0.2917}}
  >{\raggedright\arraybackslash}p{(\linewidth - 4\tabcolsep) * \real{0.3333}}
  >{\raggedright\arraybackslash}p{(\linewidth - 4\tabcolsep) * \real{0.3750}}@{}}
\toprule\noalign{}
\begin{minipage}[b]{\linewidth}\raggedright
Issue
\end{minipage} & \begin{minipage}[b]{\linewidth}\raggedright
Status
\end{minipage} & \begin{minipage}[b]{\linewidth}\raggedright
Section
\end{minipage} \\
\midrule\noalign{}
\endhead
\bottomrule\noalign{}
\endlastfoot
\(K_{7}\) topological classification & Certified metric (Paper I);
(\(b_{2}\),\(b_{3}\))=(21,77) absent from all known compact \(G_{2}\)
constructions; orthogonal TCS excluded by parity; complete geometric
construction remains open & §2 \\
Singularity structure & Not required: SM gauge group from
\(g_{2} \subset\)so(7) spectral structure (§5) & §2 \\
Formula selection rules & Consequence of metric uniqueness (see §9.4
note) & §3, §7 \\
\(\alpha\)\_s(RGE) = 3.7\% deviation & SUSY threshold sensitivity;
split-spectrum matching (§5.3) & §5 \\
\(\delta\)\_CP deviation & 11.3\% vs NuFIT 6.0 central; at edge of 1\(\sigma\) band (\(\pm 20\)\({}^{\circ}\)); see S2 Appendix F & §4.4.1, §9.1 \\
\(\Delta\)V(e-\(\mu\)) = 15.9\% deviation & Reduced to 0.75\% by
combined pipeline (\(\alpha\) = \(e^K\)) & §6 \\
Hidden \(E_{8}\) sector & Candidate for dark sector (dark matter, dark
energy); no direct observable coupling & §9.2 \\
Quantum gravity completion & Not addressed &, \\
\end{longtable}


\textbf{Note on the selection principle.} A natural objection to any
framework of this type is: why \emph{these} observables? At the mathematical level, the NK-certified metric has zero free parameters, and its 169 Chebyshev
coefficients are algebraically constrained to a five-generator lattice
Z{[}\(\pi\)\textsuperscript{2}, \(\pi\), 1, e,
\(\chi\)\textsubscript{9}\textsubscript{8}{]} / (\(b_{2} \cdot b_{3}\)).
Observables are functionals of this metric and are therefore algebraic in topological invariants. Within the declared formula grammar and structural-constant set, the observed relations are highly constrained and statistically non-generic, but a first-principles derivation of the full formula-selection rule remains open. Extended grammar analysis is consistent with this picture: adding TCS-level atoms (\(\chi\)(K3)=24,
\(b_{2}\)(\(M_{1}\))=11, \(b_{2}\)(\(M_{2}\))=10) to the search grammar
discovers simpler and more precise formulas for several observables
(m\_c/m\_s exact at 11+7/10; \(\Omega\)\_DE = 53/77 rational,
5\(\times\) more precise; \(\lambda\)\_H = 11/87, 7\(\times\) more
precise), suggesting that observables encode TCS structure more directly than the higher-level algebra currently used.

At the philosophical level, the residual question is: why \(G_{2}\)
holonomy? This is the framework's single underdetermined input: the
assertion that the compactification manifold is a compact 7-manifold of
\(G_{2}\) holonomy with \(b_{1}\)=0. This question has the same
epistemic status as ``why Lorentz invariance?'' in special relativity or
``why the Dirac equation?'' in relativistic quantum mechanics: it is
recognized as the minimal consistent mathematical structure compatible
with observed symmetries, not derived from something more fundamental.
\(G_{2}\) is the unique compact exceptional holonomy group admitting a
7-dimensional representation with the stabilizer chain \(G_{2}\)
\(\supset\) SU(3) \(\supset\) SU(2) \(\supset\) U(1) that naturally
contains the subgroups needed for the Standard Model. It is not selected
among a continuum of options; it is an isolated point in the space of
compatible mathematical structures. Every other unification framework
carries equivalent or greater philosophical inputs alongside
substantially more free parameters. GIFT makes its single input
explicit.

\textbf{Honest assessment of outliers}: Two observables deviate by
\textgreater5\%: \(\Delta\)V(e-\(\mu\)) (15.9\%, reduced to 0.75\% by
the combined wilson\_line+instanton pipeline) and \(\kappa\)(gauge)
(4.7\%, reflecting genuine \(K_{7}\) geometry at the percent level). The
\(\alpha\)\_s(RGE) deviation is 3.7\% with split-spectrum matching
(§5.3); a naive degenerate-spectrum treatment would give 12.1\%. The
\(\delta\)\_CP prediction (197\({}^{\circ}\)) is 11.3\% from NuFIT 6.0
central but at the edge of the 1\(\sigma\) uncertainty band
(177\({}^{\circ}\) \(\pm\) 20\({}^{\circ}\)). A post-hoc structural observation involving a 62/69 factor is recorded in S2 Appendix F (not part of the main framework).

\subsection{Conclusion}

We have explored a framework deriving \textbf{95 observables} from
topological invariants of a compact \(G_{2}\) manifold \(K_{7}\) with
\(E_{8} \times E_{8}\) gauge structure:

\begin{itemize}

\item
  \textbf{95 observables} organized in 4 types (33(I) + 19(II) + 21(III)
  + 22(IV)), with 66 testable against experiment
\item
  \textbf{Mean deviation \textasciitilde1.15\%} across 66 comparable
  observables; 0.73\% for Type I; 11 exact matches, 53 within 1\%
\item
  \textbf{Three new physics sections}: gauge breaking (\(E_{8}\) \(\to\)
  SM, §5), mass hierarchy (combined wilson\_line+instanton, §6),
  sensitivity analysis (r\_eff = 15.53, P(uniform) = 10\(^{-346}\),
  P(algebraic) = 10\(^{-133}\), §7)
\item
  \textbf{Lean 4 certification}: 213 conjuncts, 0 sorry, 55/95
  observables verified, 4 substantive axioms
\item
  \textbf{Statistical distinctiveness} at \textgreater{} 4.5\(\sigma\)
  among 192,349 alternatives tested
\item
  \textbf{Falsifiable predictions}: \(\delta\)\_CP = 197\({}^{\circ}\),
  N\_gen = 3, \(sin^{2} \theta\)\_W = 3/13, testable by DUNE/FCC-ee
\end{itemize}

\textbf{We do not claim this framework is correct.} It may represent
genuine geometric insight, effective approximation, or elaborate
coincidence. Only experiment can discriminate. The deeper question,
why octonionic geometry would determine particle physics parameters,
remains open.

But the empirical success of 95 observables from zero adjustable
parameters, with an overdetermination ratio of 2.13\(\times\) and
coincidence probability 10\(^{-346}\) (uniform) / 10\(^{-133}\)
(algebraic), suggests that topology and physics may be more intimately
connected than currently understood.

\textbf{The ultimate arbiter is experiment.}

\noindent

\rule{\textwidth}{0.2pt}

\section{Author's note}

This framework was developed through sustained collaboration between the
author and several AI systems, primarily Claude (Anthropic), with
contributions from GPT (OpenAI), Gemini (Google), Grok (xAI), for
specific mathematical insights. The formal verification in Lean 4,
architectural decisions, and many key derivations emerged from iterative
dialogue sessions over several months. This collaboration follows
transparent crediting approach for AI-assisted mathematical research.
Mathematical constants underlying these relationships represent timeless
logical structures that preceded human discovery. The value of any
theoretical proposal depends on mathematical coherence and empirical
accuracy, not origin. Mathematics is evaluated on results, not résumés.

\noindent

\rule{\textwidth}{0.2pt}

\section{Competing Interests}

The author declares no competing interests.

\noindent

\rule{\textwidth}{0.2pt}

\section*{References}
\addcontentsline{toc}{section}{References}

\input{gift_3.4_bib}

\noindent

\rule{\textwidth}{0.2pt}

\section*{Author's Related Works}

\noindent\emph{Code and Lean proofs:} \url{https://github.com/gift-framework/core} (Lean module under \texttt{core/Lean/}).

\begin{itemize}
\item \textbf{[A]} B. de La Fourni\`ere, ``An Explicit Approximate $\mathrm{G}_2$ Metric on a Compact 7-Manifold with Certified Torsion-Free Completion,'' Zenodo \href{https://doi.org/10.5281/zenodo.19892350}{10.5281/zenodo.19892350} (2026).
\item \textbf{[B]} B. de La Fourni\`ere, ``Spectral Geometry of the $\mathrm{G}_2$-GIFT Manifold: Betti Numbers, KK Spectrum, and Spectral Invariants,'' Zenodo \href{https://doi.org/10.5281/zenodo.19893371}{10.5281/zenodo.19893371} (2026).
\item \textbf{[C]} B. de La Fourni\`ere, ``Newton-Kantorovich Certificate for the K3 Donaldson Embedding in the $\mathrm{G}_2$-GIFT Metric,'' Zenodo \href{https://doi.org/10.5281/zenodo.19708916}{10.5281/zenodo.19708916} (2026).
\item \textbf{[D]} B. de La Fourni\`ere, ``An Explicit Closed-Form $\mathrm{G}_2$ Ansatz on a K3-Coassociative Neck with Hyperk\"ahler Rotation and Picard--Lefschetz Wirtinger Certificate,'' Zenodo \href{https://doi.org/10.5281/zenodo.20039066}{10.5281/zenodo.20039066} (2026).
\item \textbf{[essay]} B. de La Fourni\`ere, \emph{Orientation, not ontology} (companion essay), \url{https://giftheory.substack.com/p/orientation-not-ontology} (2026).
\end{itemize}

\noindent

\rule{\textwidth}{0.2pt}

\section{Appendix A: Topological Input Constants}

\begin{longtable}[]{@{}lll@{}}
\toprule\noalign{}
Symbol & Definition & Value \\
\midrule\noalign{}
\endhead
\bottomrule\noalign{}
\endlastfoot
dim(\(E_{8}\)) & Lie algebra dimension & 248 \\
rank(\(E_{8}\)) & Cartan subalgebra dimension & 8 \\
dim(\(K_{7}\)) & Manifold dimension & 7 \\
\(b_{2}\)(\(K_{7}\)) & Second Betti number & 21 \\
\(b_{3}\)(\(K_{7}\)) & Third Betti number & 77 \\
dim(\(G_{2}\)) & Holonomy group dimension & 14 \\
dim(\(J_{3}\)(O)) & Jordan algebra dimension & 27 \\
\end{longtable}


The complete set of 20 structural constants (including dim(\(E_{6}\)),
dim(\(E_{7}\)), dim(\(F_{4}\)), fund(\(E_{7}\)),
\textbar PSL(2,7)\textbar, D\_bulk, \(\alpha\)\_sum, det(g)\_den,
det(g)\_num) is tabulated in Supplement S3, §3.3.

\section{Appendix B: Derived Structural Constants}

\begin{longtable}[]{@{}lll@{}}
\toprule\noalign{}
Symbol & Formula & Value \\
\midrule\noalign{}
\endhead
\bottomrule\noalign{}
\endlastfoot
\(p_{2}\) & dim(\(G_{2}\))/dim(\(K_{7}\)) & 2 \\
Weyl & From W(\(E_{8}\)) factorization & 5 \\
N\_gen & Index theorem & 3 \\
H* & \(b_{2}\) + \(b_{3}\) + 1 & 99 \\
tau & (496 x 21)/(27 x 99) & 3472/891 \\
kappa\_T & 1/(\(b_{3}\) - dim(\(G_{2}\)) - \(p_{2}\)) & 1/61 \\
det(g) & \(p_{2}\) + 1/(\(b_{2}\) + dim(\(G_{2}\)) - N\_gen) & 65/32 \\
\end{longtable}


\section{Appendix C: Supplement Reference}

\begin{longtable}[]{@{}
  >{\raggedright\arraybackslash}p{(\linewidth - 4\tabcolsep) * \real{0.4286}}
  >{\raggedright\arraybackslash}p{(\linewidth - 4\tabcolsep) * \real{0.3214}}
  >{\raggedright\arraybackslash}p{(\linewidth - 4\tabcolsep) * \real{0.2500}}@{}}
\toprule\noalign{}
\begin{minipage}[b]{\linewidth}\raggedright
Supplement
\end{minipage} & \begin{minipage}[b]{\linewidth}\raggedright
Content
\end{minipage} & \begin{minipage}[b]{\linewidth}\raggedright
Pages
\end{minipage} \\
\midrule\noalign{}
\endhead
\bottomrule\noalign{}
\endlastfoot
S1: Foundations & \(E_{8}\), \(G_{2}\), \(K_{7}\) construction details &
27 \\
S2: Derivations & Complete proofs of 33 Type I relations &
42 \\
S3: Observable Dataset & Full 95-entry table, type classification,
sensitivity & 10 \\
\end{longtable}


\noindent

\rule{\textwidth}{0.2pt}

\emph{GIFT Framework}

\end{document}
